% =====================================================================
%  Informe de bloqueos - ATENEA2 (SRVPRORPDB01)
%  Departamento de IT - Servicios Generales de la Plana
%  Compilar con:  pdflatex informe-bloqueos-atenea2.tex   (dos pasadas)
% =====================================================================
\documentclass[11pt,a4paper]{article}

\usepackage[utf8]{inputenc}
\usepackage[T1]{fontenc}
\usepackage[spanish,es-noquoting]{babel}
\usepackage[a4paper,margin=2.4cm]{geometry}
\usepackage{lmodern}
\usepackage{booktabs}
\usepackage{longtable}
\usepackage{array}
\usepackage{xcolor}
\usepackage{colortbl}   % \cellcolor en las tablas del plan por fases
\usepackage{listings}
\usepackage{graphicx}
\usepackage{tikz}
\usepackage{pgfplots}
\pgfplotsset{compat=1.18}
\usepackage[hidelinks]{hyperref}
\usepackage{fancyhdr}
\usepackage{titlesec}

% ----------------------------- colores -------------------------------
\definecolor{tinta}{HTML}{1A1A1A}
\definecolor{acento}{HTML}{1F6FEB}
\definecolor{gris}{HTML}{6B7280}
\definecolor{grisclaro}{HTML}{F3F4F6}
\definecolor{rojo}{HTML}{B91C1C}
\definecolor{verde}{HTML}{15803D}
\definecolor{ambar}{HTML}{B45309}
\definecolor{comentario}{HTML}{6B7280}
\definecolor{cadena}{HTML}{15803D}

% --------------------------- cabecera/pie ----------------------------
\pagestyle{fancy}
\fancyhf{}
\renewcommand{\headrulewidth}{0.4pt}
\fancyhead[L]{\footnotesize\color{gris} Informe de bloqueos · ATENEA2}
\fancyhead[R]{\footnotesize\color{gris} 29 de octubre de 2026}
\fancyfoot[C]{\footnotesize\color{gris} SGP · Departamento de IT}
\fancyfoot[R]{\footnotesize\color{gris}\thepage}

\titleformat{\section}{\Large\bfseries\color{tinta}}{\thesection}{0.6em}{}
\titleformat{\subsection}{\large\bfseries\color{tinta}}{\thesubsection}{0.6em}{}

% ------------------------- estilo de codigo --------------------------
\lstdefinestyle{sql}{
  language=SQL,
  basicstyle=\ttfamily\footnotesize,
  keywordstyle=\color{acento}\bfseries,
  commentstyle=\color{comentario}\itshape,
  stringstyle=\color{cadena},
  backgroundcolor=\color{grisclaro},
  frame=leftline,
  framerule=2pt,
  rulecolor=\color{acento},
  framexleftmargin=6pt,
  xleftmargin=10pt,
  breaklines=true,
  breakatwhitespace=true,
  showstringspaces=false,
  captionpos=b,
  numbers=left,
  numberstyle=\tiny\color{gris},
  numbersep=10pt,
  tabsize=2,
  morekeywords={OVER,PARTITION,CAST,DECIMAL,DATEDIFF,SYSDATETIME,OUTER,APPLY,
                NULLIF,SUM,MAX,COUNT,DISTINCT,TOP,EXEC,CONVERT,uniqueidentifier,
                SECOND,DAY,BIGINT,INT,DB_NAME,OBJECT_NAME}
}

\lstdefinestyle{salida}{
  basicstyle=\ttfamily\scriptsize,
  backgroundcolor=\color{grisclaro},
  frame=leftline,
  framerule=2pt,
  rulecolor=\color{gris},
  framexleftmargin=6pt,
  xleftmargin=10pt,
  breaklines=true,
  showstringspaces=false,
  columns=fullflexible
}

\newcommand{\cod}[1]{\texttt{\small #1}}
\newcommand{\critico}[1]{\textcolor{rojo}{\textbf{#1}}}
\newcommand{\bien}[1]{\textcolor{verde}{\textbf{#1}}}
\newcommand{\aviso}[1]{\textcolor{ambar}{\textbf{#1}}}
\newcommand{\sombra}{\cellcolor{grisclaro}}

% =====================================================================
\begin{document}

\begin{titlepage}
\centering
\vspace*{3cm}
{\Huge\bfseries Informe de bloqueos\par}
\vspace{0.6cm}
{\LARGE ATENEA2 \textemdash{} SRVPRORPDB01\par}
\vspace{1.4cm}
{\large Análisis de bloqueos recurrentes en instancia SQL Server\par}
\vspace{0.4cm}
{\large\color{gris} Monitorizada con SolarWinds DPA 2022.4\par}
\vfill
{\large Departamento de IT\par}
{\large Servicios Generales de la Plana\par}
\vspace{0.8cm}
{\large 29 de octubre de 2026\par}
\vspace{2cm}
\end{titlepage}

\tableofcontents
\newpage

% =====================================================================
\section{Resumen}

La instancia ATENEA2 (SRVPRORPDB01) presenta bloqueos de forma
\textbf{permanente durante toda la jornada laboral}, no de forma puntual.
En los últimos siete días, el \critico{88\,\%} de las muestras registran
sesiones bloqueadas. Las otras cuatro instancias monitorizadas se mueven
entre el 0\,\% y el 2\,\%.

El desencadenante no son los bloqueos en sí. De las 60 sesiones
examinadas, \textbf{45 esperaban por paralelismo} y solo 7 por bloqueo de
objetos. Las sesiones que bloquean a las demás no esperan a nadie: esperan
turno de CPU.

\subsection*{Lo que se comprobó directamente en la instancia}

\begin{itemize}
  \item \bien{El paralelismo ya estaba bien configurado.} MAXDOP a 8 sobre
        16 núcleos en 2 nodos NUMA, y \emph{cost threshold} a 100 frente al
        valor por defecto de 5. Hipótesis inicial descartada.
  \item \bien{TempDB está bien dimensionado.} Ocho ficheros de datos, al
        99\,\% libre. Descartado como causa.
  \item \critico{140 planificaciones de SQL Agent sub-diarias activas},
        con seis jobs cada minuto y tres cada diez segundos. Dos de ellos
        —\cod{D365\_SINCRONIZAR\_OMS\_DISP\_MAX} al 118\,\% y
        \cod{D365\_SYNC\_TRANSFORMAR\_DISP\_MAX} al 96\,\%— \textbf{tardan
        más que su propio intervalo}: no paran nunca.
  \item \critico{El lote nocturno invade la jornada.}
        \cod{TAREAS NOCTURNAS STN} puede terminar a las 09:14,
        \cod{RAPPEL\_ZCALCULOSFACTURAS} a las 15:24 e \cod{IndexOptimize} a
        las 16:29. El primero atraviesa el pico de bloqueos de las 08:00; el
        último explicaría el bloqueo de esquema de 392\,s.
  \item \critico{Las estadísticas nunca están al día.}
        \cod{entEFECTO\_FacturasCobradas365} acumula 31,5 millones de
        modificaciones sobre 1,18 millones de filas en nueve horas: 26 veces
        su tamaño. \textbf{Esto explica el paralelismo excesivo} — el
        optimizador estima mal y elige planes paralelos que no hacen falta.
  \item \critico{Ninguna base de datos usa versionado de filas.} RCSI
        desactivado en las ocho revisadas: los lectores bloquean a los
        escritores. Es el mayor margen de mejora sobre el síntoma, y el
        cambio más delicado.
  \item \critico{La tabla \cod{entEMPRESA} acumula 38 horas de espera por
        bloqueos} en 27 días, y \cod{entORDENCARGA} otras 22. Tres jobs
        reescriben \cod{entEMPRESA} cada media hora en pleno horario laboral.
        Es la fuente de bloqueo concreta más importante localizada.
  \item \critico{Demostrado el paralelismo innecesario:} seis variantes del
        mismo \cod{INSERT} consumen \textbf{31.000 segundos de CPU} para
        devolver \textbf{cero filas}, y una consulta de \cod{OMS} recibe
        \textbf{506\,MB de memoria para devolver una}. Confirma que la causa
        está en las estimaciones, no en la configuración.
  \item \critico{La base de datos \cod{Cenusa} tiene 575 disparadores} sobre
        sus tablas, nueve de ellos sobre \cod{entEMPRESA}. Se ejecutan dentro
        de la transacción que los dispara, alargando directamente la
        retención de bloqueos. Un solo procedimiento consume
        \textbf{1,1\,GB de memoria por ejecución} y aprovecha el 17\,\%.
  \item \critico{Hay 25\,GB de memoria física sin usar} mientras SQL Server,
        topado en 94,6\,GB, trabaja con un \emph{page life expectancy} de
        28 segundos. El sistema operativo confirma que esa RAM está
        realmente disponible.
  \item \bien{El disco no es el cuello de botella.} Medido fichero por
        fichero, la peor latencia de lectura de la instancia es de 9\,ms y
        las escrituras de log van a 1--2\,ms. Confirmado también desde el
        hipervisor: los siete discos de fibra del host por debajo de 1,5\,ms.
  \item \bien{El hipervisor está descartado.} Con el vCenter ya registrado en
        DPA: \emph{CPU ready time} 0,09\,\%, \emph{co-stop} 0\,\%, sin
        \emph{ballooning} ni intercambio, host al 40\,\% de CPU. \textbf{El
        problema está íntegramente dentro de SQL Server}, y subir el tope de
        memoria queda desbloqueado sin reservas.
  \item \critico{Un índice ausente destaca sobre todos los demás:}
        \cod{Grupo.cmCargaConfirmadaLogs} sobre \cod{[TipoAccion]}, con 14,9
        millones de búsquedas en 27 días y un impacto estimado 56 veces
        mayor que el siguiente candidato.
\end{itemize}

\vspace{0.3cm}
\noindent\fbox{\parbox{0.97\textwidth}{\textbf{Acción recomendada en primer
lugar:} activar \cod{AUTO\_UPDATE\_STATISTICS\_ASYNC} en \cod{Grupo} y poner
estadísticas con \cod{FULLSCAN} varias veces al día en las dos tablas de
mayor rotación. Es un comando y un job, sin ventana de parada, reversible, y
ataca el mecanismo que genera el paralelismo excesivo en lugar del
síntoma.}}

% =====================================================================
\section{Qué se detectó}

Dos capturas del mismo día, separadas por minutos, con el panel de DPA como
fuente.

\begin{table}[h]
\centering
\small
\begin{tabular}{lrrrr}
\toprule
\textbf{Momento} & \textbf{Activas} & \textbf{Bloqueadas} &
\textbf{Pico bloq. 1\,h} & \textbf{Pico act. 1\,h} \\
\midrule
Primera lectura & 125 & 23 & 37 & --- \\
10:31           &  45 &  9 & 37 & 287 \\
10:42           & 202 &  9 & 37 & 287 \\
\bottomrule
\end{tabular}
\caption{Evolución durante la toma de datos}
\end{table}

En la primera lectura la cadena de bloqueo era de dos niveles:

\begin{table}[h]
\centering
\scriptsize
\begin{tabular}{llllll}
\toprule
\textbf{SPID} & \textbf{Rol} & \textbf{Bloq. por} & \textbf{Programa} &
\textbf{Espera} & \textbf{Tiempo} \\
\midrule
2300 & bloquea & --- & SQLAgent TSQL JobStep & SOS\_SCHEDULER\_YIELD & 0\,s \\
1184 & espera  & 2300 & .Net SqlClient       & LCK\_M\_SCH\_S & 200\,s \\
1375 & espera  & 2300 & SQLCMD               & LCK\_M\_SCH\_S & 392\,s \\
\bottomrule
\end{tabular}
\caption{Cadena de bloqueo, primera lectura}
\end{table}

A las 10:31 ya era de tres niveles, con el SPID 2375 bloqueado y bloqueando
a la vez. En ambos casos \textbf{la raíz de la cadena es un job de SQL
Agent}, y en ambos el bloqueador espera CPU o paralelismo, no un bloqueo
ajeno.

\subsection{Composición de las sesiones}

\begin{table}[h]
\centering
\small
\begin{tabular}{p{3cm}p{10.5cm}}
\toprule
\textbf{Dimensión} & \textbf{Reparto de las 60 sesiones examinadas} \\
\midrule
Programa       & 25 .Net SqlClient · 33 SQLAgent TSQL JobStep (8 jobs distintos) \\
Host           & 35 SRVPRORPDB01 (los propios jobs) · 20 SRVTAREAS · 5 puestos XenApp \\
Base de datos  & 20 OMS · 20 Grupo · 11 IFR · 9 Cenusa \\
Tipo de espera & 26 CXCONSUMER · 7 EXECSYNC · 7 CXPACKET · 5 CXSYNC\_PORT · 5 LCK\_M\_IS · 2 LCK\_M\_S \\
\bottomrule
\end{tabular}
\end{table}

Al menos diez sesiones tenían transacciones abiertas estando en estado
\cod{suspended}, la mayoría con dos cada una. Una transacción abierta en una
sesión suspendida retiene bloqueos sin estar avanzando.

% =====================================================================
\newpage
\section{Qué tan recurrente es}

No es un incidente. En los últimos siete días el \critico{88\,\%} de las
muestras registran sesiones bloqueadas, con una media de 4,9. El perfil
horario se repite día tras día: cae de madrugada, arranca a las 5\,h, se
dispara a las 8\,h y no baja hasta el final de la jornada.

\begin{figure}[h]
\centering
\begin{tikzpicture}
\begin{axis}[
    width=15cm, height=7cm,
    xlabel={Hora del día}, ylabel={Sesiones bloqueadas (media)},
    xmin=-0.5, xmax=23.5, ymin=0, ymax=10,
    xtick={0,3,6,9,12,15,18,21},
    xticklabels={00\,h,03\,h,06\,h,09\,h,12\,h,15\,h,18\,h,21\,h},
    ytick={0,3,6,9},
    grid=major, grid style={gray!20},
    legend pos=north east,
    legend style={draw=gray!40, font=\small},
    tick label style={font=\small, color=gris},
    label style={font=\small, color=gris},
    axis lines=left
]
\addplot[color=acento, line width=1.6pt, mark=*, mark size=1.2pt]
  coordinates {(0,1.6)(1,2.1)(2,1.7)(3,1.9)(4,2.3)(5,4.0)(6,4.3)(7,4.6)
               (8,9.0)(9,6.4)(10,5.9)(11,6.7)(12,8.1)(13,8.0)(14,6.0)
               (15,6.7)(16,7.7)(17,5.4)(18,5.1)(19,4.7)(20,4.4)(21,3.7)
               (22,4.4)(23,2.6)};
\addlegendentry{Media 7 días}
\addplot[color=gris, line width=1pt, dashed, mark=*, mark size=0.8pt]
  coordinates {(0,0.3)(1,0.3)(2,0.0)(3,0.0)(4,0.3)(5,2.0)(6,2.2)(7,2.0)
               (8,7.3)(9,8.0)(10,4.1)(11,4.2)(12,6.3)(13,5.7)(14,4.2)
               (15,5.0)(16,6.3)(17,3.5)(18,3.3)(19,2.2)(20,1.8)(21,2.0)
               (22,2.8)(23,0.8)};
\addlegendentry{Últimas 24 h}
\addplot[color=gris!70, line width=1pt, dotted] coordinates {(0,0.02)(23,0.02)};
\addlegendentry{Resto de instancias}
\end{axis}
\end{tikzpicture}
\caption{Perfil horario de sesiones bloqueadas en ATENEA2. Fuente:
histórico de \cod{Blocked Sessions}, 22--29 de octubre de 2026.}
\end{figure}

\subsection{Contraste con el resto del parque}

\begin{table}[h]
\centering
\small
\begin{tabular}{lrrr}
\toprule
\textbf{Instancia} & \textbf{Máximo} & \textbf{Media} &
\textbf{\% muestras con bloqueo} \\
\midrule
\critico{ATENEA2} & \critico{20} & \critico{4,89} & \critico{88\,\%} \\
SCALEPROD &  3 & 0,02 & 1\,\% \\
AQUILES   &  1 & 0,02 & 2\,\% \\
PROMETEO  &  1 & 0,01 & 1\,\% \\
SRVPBI    &  0 & 0,00 & 0\,\% \\
\bottomrule
\end{tabular}
\caption{Sesiones bloqueadas, 7 días}
\end{table}

ATENEA2 tiene entre 200 y 500 veces más bloqueo medio que cualquiera de las
otras cuatro. No es un problema del parque ni de la versión de SQL Server:
es de esta instancia.

\begin{quote}\small\color{gris}
El histórico se diezma a una muestra por hora en la ventana de siete días,
así que esos máximos son cotas inferiores. Midiendo a resolución completa
en una sola hora, el pico fue de 37 sesiones bloqueadas.
\end{quote}

% =====================================================================
\newpage
\section{Causa}

\subsection{El ciclo que se realimenta}

No hay una causa única. Hay \textbf{tres condiciones de fondo} que se
combinan, y un ciclo que una vez arrancado se realimenta.

\begin{figure}[h]
\centering
\begin{tikzpicture}[
  caja/.style={draw=gray!60, rounded corners=4pt, minimum width=4.4cm,
               minimum height=1.35cm, align=center, font=\footnotesize},
  raiz/.style={caja, draw=rojo!70, line width=1pt, fill=rojo!6,
               minimum height=1.6cm},
  destacada/.style={caja, draw=acento, line width=1.2pt, fill=acento!8},
  flecha/.style={-latex, draw=gray!70, line width=0.9pt},
  fina/.style={-latex, draw=rojo!50, line width=0.8pt}
]
% --- condiciones de fondo ---
\node[raiz] (est)  at (-5.2, 4.8) {\textbf{Estadísticas obsoletas}\\
    \scriptsize\color{gris} 31,5\,M de cambios sobre\\[-1pt]
    \scriptsize\color{gris} 1,18\,M de filas};
\node[raiz] (job)  at ( 0.0, 4.8) {\textbf{140 jobs sub-diarios}\\
    \scriptsize\color{gris} dos de ellos al 118\,\%\\[-1pt]
    \scriptsize\color{gris} y al 96\,\% de ocupación};
\node[raiz] (rcsi) at ( 5.2, 4.8) {\textbf{Sin versionado de filas}\\
    \scriptsize\color{gris} RCSI desactivado\\[-1pt]
    \scriptsize\color{gris} en las ocho bases};

% --- mecanismo ---
\node[destacada] (plan) at (-3.6, 2.2) {\textbf{Estimaciones erróneas}\\
    \scriptsize\color{gris} planes paralelos innecesarios};
\node[caja]      (escr) at ( 3.6, 2.2) {\textbf{Escritura constante}\\
    \scriptsize\color{gris} y lectores que bloquean};

\node[destacada] (cpu)  at (-3.6, -0.3) {\textbf{CPU saturada}\\
    \scriptsize\color{gris} 72--83\,\%; 58,7\,\% de esperas\\[-1pt]
    \scriptsize\color{gris} son de reparto paralelo};
\node[caja]      (tran) at ( 3.6, -0.3) {\textbf{Transacciones largas}\\
    \scriptsize\color{gris} retienen bloqueos más tiempo};

\node[caja]      (bloq) at ( 0.0, -2.8) {\textbf{Bloqueo permanente}\\
    \scriptsize\color{gris} 88\,\% de las muestras, pico de 9 a las 8\,h};

% --- flechas ---
\draw[fina]   (est)  -- (plan);
\draw[fina]   (job)  -- (escr);
\draw[fina]   (rcsi) -- (escr);
\draw[flecha] (plan) -- (cpu);
\draw[flecha] (escr) -- (tran);
\draw[flecha] (cpu)  -- (tran);
\draw[flecha] (cpu)  -- (bloq);
\draw[flecha] (tran) -- (bloq);
\draw[flecha] (bloq) to[out=200,in=250,looseness=1.6]
      node[below left,font=\scriptsize,color=gris,pos=0.45]{se realimenta} (cpu);
\end{tikzpicture}
\caption{Tres condiciones de fondo (en rojo) alimentan un ciclo que se
sostiene a sí mismo. El hipervisor, el disco y la configuración de
paralelismo quedaron descartados por medición.}
\end{figure}

La realimentación importa para entender por qué el cuadro no se corrige solo:
las sesiones que esperan turno de CPU retienen sus bloqueos mientras esperan,
y los usuarios que ven la aplicación colgada cierran y vuelven a entrar,
añadiendo sesiones a la misma cola.

\subsection{Las esperas lo confirman}

Acumulado de los 27 días que lleva encendida la instancia (arrancó el 2 de
octubre de 2026):

\begin{table}[h]
\centering
\small
\begin{tabular}{lrl}
\toprule
\textbf{Tipo de espera} & \textbf{\% del total} & \textbf{Qué indica} \\
\midrule
CXCONSUMER                & 27,6\,\% & paralelismo \\
CXSYNC\_PORT              & 18,8\,\% & paralelismo \\
CXPACKET                  & 12,4\,\% & paralelismo \\
SOS\_SCHEDULER\_YIELD     &  7,8\,\% & \critico{falta de CPU} \\
LATCH\_EX                 &  6,0\,\% & contención interna \\
\sombra HADR\_NOTIFICATION\_DEQUEUE & \sombra 2,8\,\% &
  \sombra \emph{espera inactiva --- no computa} \\
HTDELETE                  &  2,2\,\% & uniones hash en paralelo \\
OLEDB                     &  2,1\,\% & consultas distribuidas \\
PAGEIOLATCH\_SH           &  2,1\,\% & \critico{lecturas a disco por falta de memoria} \\
\bottomrule
\end{tabular}
\end{table}

\paragraph{Corrección sobre una lectura anterior de esta tabla.} En una
versión previa de este informe se interpretó
\cod{HADR\_NOTIFICATION\_DEQUEUE} como señal de carga del grupo de
disponibilidad. \textbf{Es un error}: esa espera, como
\cod{HADR\_WORK\_QUEUE}, \cod{HADR\_TIMER\_TASK} y
\cod{HADR\_FILE\-STREAM\_IOMGR\_IOCOMPLETION}, corresponde a hilos de fondo
ociosos esperando trabajo. Aparece con cifras enormes en cualquier instancia
con AG y debe excluirse del análisis, igual que \cod{SLEEP\_*} o
\cod{XE\_TIMER\_EVENT}. No indica nada sobre el rendimiento.

El paralelismo suma el \critico{58,7\,\%} de todo el tiempo de espera
\emph{pese a estar bien configurado}. El problema no es que se reparta
trabajo que no debería repartirse, sino que \textbf{se reparte trabajo que no
hacía falta repartir} — en la sección \ref{sec:prueba} se demuestra con
consultas concretas.

\subsection{Quince jobs a la vez, ocho en el mismo segundo}

Este fue el primer indicio. Resultó ser la punta de algo bastante mayor, que
se detalla en los dos apartados siguientes.

\begin{table}[h]
\centering
\small
\begin{tabular}{lr}
\toprule
\textbf{Hora de arranque} & \textbf{Jobs que arrancan} \\
\midrule
\critico{11:10:02} & \critico{8} \\
11:00:01 & 3 \\
11:09:08 & 1 \\
11:08:06 & 1 \\
11:02:36 & 1 \\
11:02:06 & 1 \\
\bottomrule
\end{tabular}
\caption{Jobs de SQL Agent en ejecución simultánea}
\end{table}

Los ocho que coinciden en el mismo segundo son, en su mayoría,
sincronizaciones con Dynamics 365 y con SCALE:

\begin{itemize}\itemsep1pt
  \item \cod{BLOQUEAR EMPRESAS CON IMPAGOS}
  \item \cod{D365\_SINCRONIZAR\_OMS\_DISP\_MAX}
  \item \cod{D365\_SYNC\_MANTENER\_PEDIDO\_GRUPO\_DESDED365}
  \item \cod{STN\_ERP\_StnCRMSOLICITUD\_AprobacionAuto}
  \item \cod{STN\_MARCAR\_LINEAS\_PREPARADO\_TRASLADADO}
  \item \cod{STN\_SINCRONIZASTOCKRUMBO}
  \item \cod{stnSCALE\_COPIA\_SHIPPING\_LOAD}
  \item \cod{stnSCALE\_INT\_ENT\_expediciones}
\end{itemize}

\textbf{Ocho procesos pesados disparados en el mismo segundo, sobre 16
núcleos, cada diez minutos.} Encaja con todo lo demás: la ola de
paralelismo, el \cod{SOS\_SCHEDULER\_YIELD}, el desplome del \emph{page life
expectancy} y las cadenas de bloqueo que aparecen y se deshacen en segundos.

\subsection{Pero no son quince jobs: son ciento cuarenta}

La foto anterior era un instante. Al consultar \cod{msdb.dbo.sysschedules}
aparece el cuadro completo: \critico{140 planificaciones sub-diarias activas},
sin contar las que corren una vez al día.

\begin{table}[h]
\centering
\small
\begin{tabular}{lrlr}
\toprule
\textbf{Cada} & \textbf{Jobs} & \textbf{Cada} & \textbf{Jobs} \\
\midrule
10\,s  & 3  & 5\,min  & 13 \\
30\,s  & 5  & 10\,min & 6 \\
75--100\,s & 5 & 15\,min & 9 \\
1\,min & 6  & 30\,min & 10 \\
2\,min & 7  & 1\,h    & 24 \\
3\,min & 10 & 2--4\,h & 15 \\
4\,min & 4  & 6--12\,h & 17 \\
\bottomrule
\end{tabular}
\caption{Planificaciones sub-diarias activas en SRVPRORPDB01}
\end{table}

Seis jobs cada minuto, siete cada dos, diez cada tres, trece cada cinco. Tres
jobs \emph{cada diez segundos}. En siete días, solo los veinte más frecuentes
suman en torno a \textbf{133.000 ejecuciones}, unas 19.000 al día.

\subsubsection*{Y algunos tardan más que su propio intervalo}

Es el dato que convierte el cuadro en un problema estructural. Comparando la
duración media de cada job con el intervalo al que está programado:

\begin{table}[h]
\centering
\small
\begin{tabular}{lrrrr}
\toprule
\textbf{Job} & \textbf{Programado} & \textbf{Media} & \textbf{Máximo} &
\textbf{Ocupación} \\
\midrule
D365\_SINCRONIZAR\_OMS\_DISP\_MAX  & 120\,s & 142\,s & 671\,s  & \critico{118\,\%} \\
D365\_SYNC\_TRANSFORMAR\_DISP\_MAX & 180\,s & 173\,s & 689\,s  & \critico{96\,\%} \\
D365\_ENVIO\_RECEPCIONES           &  60\,s &  47\,s & 913\,s  & \critico{78\,\%} \\
stnSCALE\_COPIA\_SHIPPING\_LOAD    &  10\,s &   7\,s & 476\,s  & \critico{70\,\%} \\
D365\_SINCRONIZAR\_DATOS\_COMPRAS  & 300\,s & 200\,s & 707\,s  & \aviso{67\,\%} \\
D365\_SYNC\_MANTENER\_RESERVAS     & 240\,s &  98\,s & 667\,s  & 41\,\% \\
D365\_SYNC\_CREAR\_PEDIDO\_GRUPO   & 240\,s &  97\,s & 1.737\,s & 40\,\% \\
\bottomrule
\end{tabular}
\end{table}

\cod{D365\_SINCRONIZAR\_OMS\_DISP\_MAX} está programado cada dos minutos y
tarda de media dos minutos y veintidós segundos. El «cada dos minutos» es
ficción: ese job no para nunca. SQL Agent no solapa una instancia consigo
misma, así que lo que ocurre es que arranca de nuevo en cuanto termina, de
forma continua, todo el día. \cod{D365\_SYNC\_TRANSFORMAR\_DISP\_MAX}, al
96\,\%, está en la misma situación de hecho.

\subsubsection*{Los máximos delatan quién sufre el bloqueo}

Hay jobs cuya duración máxima no guarda ninguna relación con su media:

\begin{table}[h]
\centering
\small
\begin{tabular}{lrrr}
\toprule
\textbf{Job} & \textbf{Media} & \textbf{Máximo} & \textbf{Factor} \\
\midrule
Interface APP WEB                 & 16\,s & 4.590\,s (76\,min) & \critico{287$\times$} \\
stnSCALE\_INT\_ENT\_PRODUCCION    &  5\,s & 1.781\,s (30\,min) & \critico{356$\times$} \\
STNAUTO\_AUTOMATIZAR\_ALBARANEADO &  8\,s & 1.161\,s (19\,min) & \critico{145$\times$} \\
STN\_MARCAR\_LINEAS\_PREPARADO    & 230\,s & 2.336\,s (39\,min) & 10$\times$ \\
\bottomrule
\end{tabular}
\end{table}

Un proceso que normalmente tarda cinco segundos y un día tarda treinta
minutos no se ha vuelto lento: se ha quedado esperando. Estos máximos no son
causa, son \textbf{consecuencia}, y demuestran que el bloqueo no solo molesta
a los usuarios, también atasca la propia automatización.

\subsubsection*{El lote nocturno invade la jornada}

Y por último, varios procesos nocturnos no terminan por la noche:

\begin{table}[h]
\centering
\small
\begin{tabular}{llrl}
\toprule
\textbf{Job} & \textbf{Arranca} & \textbf{Máx.} & \textbf{Termina (peor caso)} \\
\midrule
RAPPEL\_ZCALCULOSFACTURAS & 04:00 & 11\,h 24\,m & \critico{15:24} \\
IndexOptimize USER\_DATABASES Weekend & 02:30 & 13\,h 59\,m & \critico{16:29} \\
TAREAS NOCTURNAS STN & 01:00 & 8\,h 14\,m & \critico{09:14} \\
STN\_MANTENER\_ARTICULOS\_CONSUMOS & --- & 3\,h 09\,m & --- \\
BAJA\_RIESGOS & 04:00 & 2\,h 16\,m & 06:16 \\
ACTUALIZAR FOTO CARTERA ARTICULO & 05:30 & 1\,h 24\,m & 06:54 \\
\bottomrule
\end{tabular}
\end{table}

\cod{TAREAS NOCTURNAS STN} puede terminar a las 09:14, lo que significa que
\textbf{atraviesa de lleno el pico de bloqueos de las 08:00} documentado en el
apartado de recurrencia.

El historial de \cod{IndexOptimize} de los últimos ocho días confirma que el
desbordamiento es real y recurrente:

\begin{table}[h]
\centering
\small
\begin{tabular}{llrl}
\toprule
\textbf{Fecha} & \textbf{Empieza} & \textbf{Duración} & \textbf{Termina} \\
\midrule
24/10 (sáb, \emph{Weekend}) & 02:30 & 13\,h 59\,m & \critico{16:29} \\
26/10 (dom, \emph{Daily})   & 02:30 &  6\,h 29\,m & \critico{08:59} \\
22/10 & 02:30 & 1\,h 25\,m & 03:55 \\
21, 23, 27, 28, 29/10 & 02:30 & $\approx$57\,min & $\approx$03:27 \\
25/10 & 02:30 & \multicolumn{2}{l}{\aviso{registro corrupto en el historial}} \\
\bottomrule
\end{tabular}
\end{table}

En días normales termina a las 03:27, pero el sábado llegó a las 16:29 y el
domingo a las 08:59. Son fines de semana, pero demuestra que la ventana
nocturna no contiene al proceso. (La fila del 25/10 devuelve una duración
negativa y una fecha de fin de 2015: un registro mal escrito en
\cod{sysjobhistory}, probablemente de una ejecución interrumpida.)

\paragraph{Corrección: \cod{IndexOptimize} no explica el bloqueo de esquema.}
En una versión anterior de este informe se propuso que la reconstrucción de
índices causaba el bloqueo \cod{Sch-M} de 392 segundos observado el 29 de
octubre por la mañana. \textbf{Los horarios lo desmienten}: ese día
\cod{IndexOptimize} se ejecutó de 02:30 a 03:26:58, siete horas antes del
incidente. La hipótesis queda descartada.

La explicación más probable ahora es otra: un \cod{TRUNCATE TABLE} sobre
alguna de las tablas de paso que se vacían y recargan continuamente
—\,\cod{entEFECTO\_FacturasCobradas365} es la candidata evidente, con sus
31,5 millones de modificaciones\,—, porque \cod{TRUNCATE} también toma un
bloqueo \cod{Sch-M}. \textbf{Es una hipótesis, no un hecho}: confirmarla
requiere capturar el \emph{blocked process report} o una sesión de
\emph{Extended Events}, que no se ha hecho.

\subsubsection*{Nota: el equipo ya conocía esto}

Entre los jobs activos hay dos cuyos nombres son, en sí mismos, un hallazgo:

\begin{itemize}\itemsep1pt
  \item \cod{REVISAR BLOQUEOS DE 7:30 a 8:30} — con tres planificaciones, una
        de ellas cada 2 minutos entre las 06:00 y las 09:00.
  \item \cod{CONTROL JOBS SINFIN} — cada 15 minutos, todo el día.
\end{itemize}

Alguien del equipo ya había identificado tanto la ventana de bloqueos de las
08:00 como la tendencia de los jobs a quedarse colgados, y construyó
vigilancia para ambas cosas. Este informe no descubre el problema: lo
cuantifica.

\subsection{Las estadísticas están permanentemente obsoletas}

Este es, probablemente, el mecanismo que explica el paralelismo excesivo, y no
se había mirado hasta ahora.

\begin{table}[h]
\centering
\small
\begin{tabular}{lrrrr}
\toprule
\textbf{Tabla (BD \cod{Grupo})} & \textbf{Filas} & \textbf{Modificaciones} &
\textbf{Churn} & \textbf{Muestra} \\
\midrule
entEFECTO\_FacturasCobradas365 & 1.184.877 & \critico{31.476.588} &
  \critico{26,6$\times$} & \critico{11,8\,\%} \\
STN\_ERP\_ANALISIS\_VENTAS\_PEDIDOS & 6.460.673 & \critico{12.940.889} &
  \critico{2,0$\times$} & \critico{1,4\,\%} \\
\bottomrule
\end{tabular}
\caption{Estadísticas actualizadas esta misma madrugada (02:41 y 02:48)}
\end{table}

Las estadísticas de estas tablas se actualizaron \emph{hoy}, de madrugada. No
son antiguas por fecha —\,\cod{dias} = 0\,—. El problema es el volumen de
cambio acumulado desde entonces: en menos de nueve horas,
\cod{entEFECTO\_FacturasCobradas365} ha registrado \textbf{31,5 millones de
modificaciones sobre una tabla de 1,18 millones de filas}. Veintiséis veces su
propio tamaño. Es una tabla de paso que se vacía y recarga continuamente, lo
que concuerda con los jobs \cod{D365 CARGAR FACTURAS COBRADAS} (cada hora),
\cod{ACTUALIZAR\_FECHACOBRO\_EFECTOS\_DESDE\_CONCILIADOS} (cada 30 minutos) y
\cod{Limpiar Fechas Efectos} (cada hora).

El umbral de actualización automática en el nivel de compatibilidad 160 es
aproximadamente $\sqrt{1000 \times filas}$, unas 34.400 modificaciones para
esta tabla. Se cruza en cuestión de minutos. En la práctica, \textbf{el
optimizador nunca trabaja con estadísticas representativas de esa tabla}.

Y la calidad del muestreo agrava el efecto: 11,8\,\% en una tabla y
\textbf{1,4\,\%} en una de 6,4 millones de filas. Estimar la cardinalidad de
6,4 millones de filas a partir del 1,4\,\% de la muestra, con el doble de su
tamaño en cambios pendientes, produce estimaciones muy alejadas de la
realidad.

\paragraph{Y esto cierra el círculo del paralelismo.} Una estimación de filas
inflada eleva el coste estimado del plan. El \emph{cost threshold for
parallelism} está en 100 —correctamente—, pero un coste \emph{estimado} a
partir de datos erróneos lo supera con facilidad. El optimizador elige
entonces un plan paralelo, con su concesión de memoria dimensionada sobre la
misma estimación errónea. De ahí los 58,7\,\% de esperas en
\cod{CXCONSUMER}, \cod{CXSYNC\_PORT} y \cod{CXPACKET}: no es que el
paralelismo esté mal configurado —ya se comprobó que no—, es que \textbf{se
está eligiendo paralelismo para consultas que no lo necesitan}, porque el
optimizador cree que van a devolver mucho más de lo que devuelven.

\paragraph{Agravante: las actualizaciones son sincrónicas.}
\cod{is\_auto\_update\_stats\_async\_on} está en \cod{false} en las ocho bases
de datos revisadas. Cuando se cruza el umbral, la consulta que lo cruza
\textbf{espera} a que se recalculen las estadísticas antes de compilar. Sobre
una tabla de 6,4 millones de filas, en horario laboral, con 200 sesiones
activas, esa espera se la come el usuario.

\subsection{La prueba: planes paralelos que no devuelven nada}
\label{sec:prueba}

Hasta aquí, la cadena «estadísticas malas $\rightarrow$ estimaciones infladas
$\rightarrow$ paralelismo innecesario» era una inferencia coherente. Estas
cifras la convierten en un hecho. La huella de una estimación inflada es doble:
el optimizador \textbf{concede memoria} para las filas que cree que va a
manejar, y \textbf{reparte la consulta entre varios núcleos} por el mismo
motivo. Si luego devuelve cero filas, la estimación era errónea.

\begin{table}[h]
\centering
\small
\begin{tabular}{lrrrrr}
\toprule
\textbf{Sentencia (BD \cod{Grupo})} & \textbf{Ejec.} & \textbf{CPU} &
\textbf{Duración} & \textbf{Núcleos} & \textbf{Filas} \\
\midrule
\cod{INSERT StnERPArticuloLoteLinea} (×6) & 2.147 & \critico{31.000\,s} &
  12.166\,s & 2,5 & \critico{0} \\
\cod{UPDATE stnERPORDENPREPARACION} & 1.122 & 1.340\,s & 376\,s &
  \critico{3,6} & \critico{0} \\
\cod{UPDATE stnERPORDENPREPARACION} (2) & 1.122 & 919\,s & 328\,s & 2,8 &
  \critico{0} \\
\cod{DELETE D365Eventos..STNAppSuite\ldots} & 591 & 1.993\,s & 898\,s & 2,2 & 3 \\
\cod{SELECT OMS..InventTransOrigin\ldots} & 164 & 581\,s & 291\,s & 2,0 &
  \critico{1} \\
\bottomrule
\end{tabular}
\caption{«Núcleos» = \cod{total\_worker\_time / total\_elapsed\_time}: por
encima de 1 significa plan paralelo. «Filas» es la media devuelta por
ejecución.}
\end{table}

Tres cosas saltan a la vista.

\paragraph{Ocho horas y media de CPU para no devolver ninguna fila.} Las seis
variantes de \cod{INSERT INTO StnERPArticuloLoteLinea} suman
\textbf{31.000 segundos de CPU} —\,más de ocho horas\,— en apenas 2.147
ejecuciones. Son 14 segundos de CPU por ejecución, repartidos entre 2,5
núcleos, con 71\,MB de memoria concedida cada vez y cero filas de resultado.
Es, con diferencia, el mayor consumidor paralelo de la instancia.

\paragraph{506\,MB de memoria para devolver una fila.} La consulta
\cod{SELECT a.SourceKey FROM OMS..InventTransOriginBiEntities} recibe
\textbf{506\,MB de concesión por ejecución} y devuelve \textbf{una fila}. No
hay forma más clara de enseñar una estimación de cardinalidad equivocada: el
optimizador cree que va a mover cientos de miles de filas y mueve una. Esa
concesión se resta del \emph{buffer pool} mientras dura la consulta, lo que
conecta directamente con el \emph{page life expectancy} de 28 segundos.

\paragraph{Seis copias de la misma sentencia.} Que el mismo \cod{INSERT}
aparezca seis veces con estadísticas propias indica que \textbf{no está
parametrizado}: cada variante literal genera su propio plan, ocupa su propia
entrada en la caché y se optimiza por separado. Parametrizar esa sentencia
reduciría a una sexta parte el trabajo de compilación y permitiría que un
único plan se reutilizara.

\paragraph{Qué cambia esto en las recomendaciones.} Nada del plan por fases,
pero sí lo refuerza: las acciones 1 y 2 —estadísticas asíncronas y
\cod{FULLSCAN} frecuente— atacan exactamente el mecanismo que estas cifras
demuestran. Y añade un objetivo nominal para la acción 10, que hasta ahora
hablaba de «las consultas del \emph{Tuning}» sin concretar: empezar por
\cod{StnERPArticuloLoteLinea} y por la consulta de
\cod{InventTransOriginBiEntities}.

\subsection{Escalado de bloqueos: dónde duele de verdad}

\cod{sys.dm\_db\_index\_operational\_stats} mide, por tabla, cuántas veces
SQL Server convirtió bloqueos de fila en bloqueos de tabla y cuánto tiempo se
pasó esperando por bloqueos. Los 27 días acumulados dan esto:

\begin{table}[h]
\centering
\small
\begin{tabular}{lrrrr}
\toprule
\textbf{Tabla (BD \cod{Grupo})} & \textbf{Escalados} & \textbf{Espera fila} &
\textbf{Espera pág.} & \textbf{Total} \\
\midrule
entEMPRESA        &  4.066 & 83.106\,s & 55.258\,s & \critico{138.364\,s (38\,h)} \\
entORDENCARGA     & 16.398 & 54.744\,s & 25.453\,s & \critico{80.197\,s (22\,h)} \\
entALMACEN\_LOTE  &  1.840 &  5.367\,s & 20.023\,s & \critico{25.390\,s (7\,h)} \\
entPEDIDO         & 18.464 &  5.556\,s & 14.226\,s & 19.782\,s \\
entPEDIDO\_DETALLE& 14.044 &  6.938\,s &  7.033\,s & 13.971\,s \\
entMOVIMIENTO     & 11.281 &  6.232\,s &  2.845\,s &  9.077\,s \\
entOC\_DETALLE    &  2.761 &  2.211\,s &  6.778\,s &  8.989\,s \\
T\_Palets         & \aviso{193.154} & 551\,s & 4.259\,s & 4.810\,s \\
APP\_PA\_StockArticulo & \aviso{68.447} & 143\,s & 0\,s & 143\,s \\
\bottomrule
\end{tabular}
\end{table}

\paragraph{El número de escalados engaña; el tiempo de espera no.}
\cod{T\_Palets} escala 193.154 veces y apenas genera espera.
\cod{entEMPRESA} escala solo 4.066 veces y acumula \textbf{38 horas de espera
por bloqueos}. Lo que importa no es cuántas veces se escala, sino sobre qué
tabla y con cuánta concurrencia alrededor.

\paragraph{Y \cod{entEMPRESA} tiene nombre y apellidos en la lista de jobs.}
Tres de los procesos recurrentes trabajan sobre esa tabla en pleno horario
laboral:

\begin{itemize}\itemsep1pt
  \item \cod{BLOQUEAR EMPRESAS CON IMPAGOS} — cada hora, de 07:10 a 19:10
  \item \cod{BLOQUEAR EMPRESAS CON IMPAGOS MAIL} — cada hora, de 07:15 a 18:15
  \item \cod{DesbloquearEmpresasSaldoNegativo} — cada 30 minutos, de 08:00 a 19:00
\end{itemize}

Una tabla maestra de empresas, que casi cualquier consulta del ERP necesita
leer, siendo reescrita cada media hora durante toda la jornada, sin versionado
de filas. Las 38 horas de espera acumulada tienen una explicación sencilla, y
\textbf{es probablemente la fuente concreta de bloqueo más importante de la
instancia}.

\paragraph{Esto refuerza la acción 9.} Con \cod{READ\_COMMITTED\_SNAPSHOT}
activo, los lectores de \cod{entEMPRESA} dejarían de tomar bloqueos
compartidos y la mayor parte de esas 38 horas desaparecería. Y para las tres
tablas de cabecera existe además una medida intermedia, más quirúrgica y
reversible en segundos:
\cod{ALTER TABLE \ldots SET (LOCK\_ESCALATION = DISABLE)}, que impide la
conversión a bloqueo de tabla a costa de más memoria de bloqueos. Conviene
probarla antes que RCSI, porque el riesgo es mucho menor.

\subsection{La base de datos \cod{Cenusa}: 575 disparadores}

Esta base de datos apareció por casualidad, al ordenar las consultas por
memoria concedida. No estaba en el foco inicial —el análisis se había
centrado en \cod{Grupo}, que es donde se veían las cadenas de bloqueo— pero
vive en la misma instancia, ocupa 46\,GB y resultó esconder lo más llamativo
de toda la investigación.

\begin{table}[h]
\centering
\small
\begin{tabular}{lr}
\toprule
\textbf{Característica de \cod{Cenusa}} & \textbf{Valor} \\
\midrule
Tamaño                              & 47.270\,MB (46\,GB) \\
Nivel de compatibilidad             & 160 \\
\cod{READ\_COMMITTED\_SNAPSHOT}     & \critico{false} \\
\cod{AUTO\_UPDATE\_STATISTICS\_ASYNC} & \critico{false} \\
\textbf{Disparadores sobre tablas}  & \critico{575} \\
De ellos, deshabilitados            & 4 \\
\bottomrule
\end{tabular}
\end{table}

\textbf{Quinientos setenta y cinco disparadores}, de los cuales solo cuatro
están desactivados. Concentrados, además, en las tablas más transitadas:

\begin{table}[h]
\centering
\small
\begin{tabular}{lrr}
\toprule
\textbf{Tabla} & \textbf{Disparadores} & \textbf{Modificaciones (27 días)} \\
\midrule
entEMPRESA           & \critico{9} &  8.759.807 \\
entALBARAN           & \critico{8} & 26.778.967 \\
entMOVIMIENTO        & \critico{8} &  5.007.080 \\
entPEDIDO            & \critico{8} & 24.107.314 \\
entPEDIDO\_DETALLE   & \critico{8} & \critico{37.738.378} \\
entARTICULO          & 6 & 11.772.885 \\
entOC\_DETALLE       & 6 & --- \\
entARTICULO\_UBICACION & 5 & --- \\
\bottomrule
\end{tabular}
\end{table}

\paragraph{Por qué esto importa para los bloqueos.} Un disparador
\textbf{se ejecuta dentro de la transacción de quien lo dispara}. No es
trabajo que ocurra «después»: es tiempo que la sentencia original pasa
reteniendo sus bloqueos. Cada \cod{UPDATE} sobre \cod{entPEDIDO\_DETALLE}
arrastra ocho disparadores antes de poder confirmar. Con 37,7 millones de
modificaciones sobre esa tabla, la amplificación es enorme.

Las cifras del coste lo confirman:

\begin{table}[h]
\centering
\small
\begin{tabular}{lrrrr}
\toprule
\textbf{Objeto} & \textbf{Ejecuciones} & \textbf{CPU} &
\textbf{Mem. concedida} & \textbf{Aprov.} \\
\midrule
\cod{ACTUALIZA\_EMPRESAS\_DIRECCIONES\ldots} & 336 & \critico{4.828\,s} &
  \critico{1.171\,MB} & \critico{17,4\,\%} \\
\cod{tU\_entALBARAN} & 21.365 & 2.927\,s & --- & --- \\
\cod{tU\_entALBARAN} & \critico{1.700.015} & 1.851\,s & 2\,MB & 40,6\,\% \\
\cod{tU\_entMOVIMIENTO} (×6) & 212.512 c/u & \critico{10.268\,s} & 1\,MB &
  0--51\,\% \\
\bottomrule
\end{tabular}
\end{table}

\paragraph{Un gigabyte de memoria para aprovechar el 17\,\%.} El
procedimiento \cod{ACTUALIZA\_EMPRESAS\_DIRECCIONES\_FICHAS\_POR\_PROVEEDOR}
recibe \textbf{1.171\,MB de concesión en cada una de sus 336 ejecuciones} y
usa el 17,4\,\%. Además lee \textbf{500.544 páginas por ejecución}, unos
3,8\,GB de datos recorridos cada vez. Es, por amplio margen, la sentencia
más cara localizada en toda la instancia, y supera incluso al caso de los
506\,MB citado antes.

\paragraph{Un disparador con 31.000 lecturas por disparo.} La primera entrada
de \cod{tU\_entALBARAN} se ejecuta 21.365 veces leyendo \textbf{31.135 páginas
cada vez}: unos 665 millones de lecturas lógicas acumuladas, solo por ese
disparador. Y la segunda entrada se dispara \textbf{1,7 millones de veces}.

\paragraph{Las seis caras del mismo disparador.}
\cod{tU\_entMOVIMIENTO} aparece seis veces porque son seis sentencias
distintas dentro del mismo cuerpo, cada una ejecutada 212.512 veces. Entre
las seis suman \textbf{10.268 segundos de CPU}, casi tres horas, y varias
aprovechan el 0\,\% de la memoria que se les concede.

\paragraph{Qué hacer con esto.} No es una acción de una tarde, y conviene
decirlo sin adornos: \textbf{575 disparadores son una decisión de
arquitectura, no un parámetro}. Revisarlos uno a uno es trabajo de
desarrollo, de meses. Pero hay dos pasos acotados y de alto rendimiento:

\begin{enumerate}\itemsep2pt
  \item \textbf{Revisar
        \cod{ACTUALIZA\_EMPRESAS\_DIRECCIONES\_FICHAS\_POR\_PROVEEDOR}.} Una
        sola sentencia, 336 ejecuciones, 1,1\,GB de concesión y 3,8\,GB
        leídos por vez. Es el mejor retorno por hora invertida de todo el
        informe.
  \item \textbf{Revisar los disparadores de \cod{entMOVIMIENTO} y
        \cod{entALBARAN}}, que son los que más CPU consumen y los que más
        alargan las transacciones ajenas.
\end{enumerate}

\paragraph{Y la tabla que lo explica todo: \cod{AUDITORIA2}.} Revisadas las
estadísticas de \cod{Cenusa}, hay una entrada que no deja lugar a dudas:

\begin{table}[h]
\centering
\small
\begin{tabular}{lrrr}
\toprule
\textbf{Tabla} & \textbf{Filas} & \textbf{Modificaciones} & \textbf{Muestra} \\
\midrule
AUDITORIA2       & \critico{29.708.129} & \critico{89.149.199} & \critico{1,0\,\%} \\
entALBARAN\_IVA  & 20.384.782 &  81.387 & \critico{1,1\,\%} \\
entPEDIDO        &    920.167 & 167.803 & \aviso{3,7--4,4\,\%} \\
\bottomrule
\end{tabular}
\end{table}

Una tabla de auditoría de \textbf{29,7 millones de filas} que acumula
\textbf{89 millones de modificaciones} —tres veces su propio tamaño— desde
que se actualizaron sus estadísticas esta madrugada. Y muestreada al
\textbf{1\,\%}.

Esto cierra el círculo de \cod{Cenusa}: los 575 disparadores escriben en
\cod{AUDITORIA2} en cada modificación. Cada \cod{UPDATE} de la aplicación
arrastra sus ocho o nueve disparadores, que a su vez escriben auditoría,
dentro de la misma transacción y reteniendo los mismos bloqueos. El coste no
es solo de CPU: es de duración de transacción, que es exactamente lo que
genera el bloqueo.

El muestreo al 1\,\% en tablas de 20 y 30 millones de filas merece revisión
por sí mismo. Es el valor que aplica \cod{IndexOptimize} por defecto en
tablas grandes, y sobre esos volúmenes produce histogramas de poca utilidad.

\paragraph{Lo que sí está bien en \cod{Cenusa}.} Los índices ausentes
\textbf{no son un problema aquí}. El mejor candidato
(\cod{entMOVIMIENTO} sobre \cod{[atrIdTipoPedido]}) tiene un impacto estimado
de 31 millones, frente a los 45.515 millones del candidato de cabeza en
\cod{Grupo}: unas 1.500 veces menos. En \cod{Cenusa} el problema son los
disparadores y el muestreo de estadísticas, no la falta de índices.

\paragraph{Y una advertencia de alcance.} Este apartado abre un frente nuevo
que \textbf{no estaba en el encargo} y que no se ha analizado con la misma
profundidad que el resto. Las cifras son sólidas, pero las conclusiones sobre
qué hacer con los disparadores requieren conocimiento funcional de la
aplicación que este informe no tiene.

\subsection{Ninguna base de datos usa versionado de filas}

\begin{table}[h]
\centering
\small
\begin{tabular}{lccc}
\toprule
\textbf{Base de datos} & \textbf{Compat.} & \textbf{RCSI} &
\textbf{Snapshot} \\
\midrule
Grupo         & 160 & \critico{false} & \critico{OFF} \\
OMS           & 160 & \critico{false} & \critico{OFF} \\
IFR           & 160 & \critico{false} & \critico{OFF} \\
ContaSTN      & 160 & \critico{false} & \critico{OFF} \\
ContaAlaplana & 160 & \critico{false} & \critico{OFF} \\
ContaVitaCer  & 160 & \critico{false} & \critico{OFF} \\
D365Eventos   & 160 & \critico{false} & \critico{OFF} \\
Informatica   & 160 & \critico{false} & \critico{OFF} \\
\bottomrule
\end{tabular}
\end{table}

\cod{READ\_COMMITTED\_SNAPSHOT} está desactivado en las ocho, y el aislamiento
\emph{snapshot} también. Significa que toda lectura toma bloqueos
compartidos: \textbf{los lectores bloquean a los escritores y viceversa}. En
un ERP con 200 sesiones concurrentes de usuario leyendo mientras 140 jobs
escriben, eso es exactamente la receta del cuadro observado.

Activar RCSI en \cod{Grupo} haría que los lectores usaran versiones de fila en
lugar de bloqueos compartidos, y la mayor parte del bloqueo lector-escritor
desaparecería. Es el cambio con mayor efecto potencial sobre el síntoma
concreto que dispara la alarma. Tiene dos contrapartidas: consume
\emph{version store} en TempDB —\,que está al 99\,\% libre con 432\,GB, así
que hay margen de sobra\,— y cambia la semántica de lectura, por lo que exige
validación con la aplicación y una ventana con acceso exclusivo para
activarlo.

Lo positivo: el nivel de compatibilidad es 160 (SQL Server 2022) en todas, y
\cod{LEGACY\_CARDINALITY\_ESTIMATION} está en \cod{false}. El estimador
moderno está en uso; no hay nada que arreglar ahí.

\subsection{La CPU: ocupada, y es SQL Server quien la ocupa}

\begin{table}[h]
\centering
\small
\begin{tabular}{lr}
\toprule
\textbf{Medida (10 minutos, 11:35--11:44)} & \textbf{Valor} \\
\midrule
CPU de SQL Server, media   & \aviso{72\,\%} \\
CPU de SQL Server, máximo  & \critico{83\,\%} \\
CPU de otros procesos      & \bien{3--6\,\%} \\
CPU libre, mínimo          & \critico{12\,\%} \\
Tareas en cola (\emph{runnable})  & 5 \\
Trabajos sin hilo asignado & \bien{0} \\
\emph{Load factor} medio / máximo & 13 / 19 \\
\bottomrule
\end{tabular}
\end{table}

Dos lecturas, una buena y una mala. La buena: no hay ningún proceso ajeno
comiéndose la máquina —antivirus, agentes de respaldo, otra instancia—. SQL
Server es el único consumidor relevante, y
\cod{work\_queue\_count} = 0 descarta el agotamiento del grupo de hilos. La
mala: 72\,\% de media con picos del 83\,\% y solo un 12\,\% de margen, con 5
tareas esperando turno de CPU, es una máquina trabajando al límite. Esto
respalda el \cod{SOS\_SCHEDULER\_YIELD} del 7,8\,\% y justifica plantear la
ampliación de vCPU —\emph{después} de corregir las causas, porque buena parte
de esa CPU se está gastando en paralelismo innecesario.

\subsection{Y hay memoria sin usar}

\begin{table}[h]
\centering
\small
\begin{tabular}{lr}
\toprule
\textbf{Medida} & \textbf{Valor} \\
\midrule
RAM física del servidor        & 131.071\,MB (128\,GB) \\
Tope de SQL Server             & 96.896\,MB (94,6\,GB) \\
Lo que SQL Server está usando  & 96.917\,MB — su tope completo \\
RAM física libre               & \critico{25.895\,MB (25\,GB)} \\
Page life expectancy           & \critico{28\,s} \\
\bottomrule
\end{tabular}
\end{table}

Dejar 33\,GB al sistema operativo en un servidor de 128\,GB es conservador.
Lo habitual sería reservar entre 16 y 20\,GB, lo que permitiría subir el
tope de SQL Server en torno a 12--15\,GB.

Y el propio sistema operativo confirma que esa RAM está realmente libre:
\cod{sys.dm\_os\_sys\_memory} devuelve \cod{system\_memory\_state\_desc} =
\emph{«Available physical memory is high»}, y
\cod{sys.dm\_os\_process\_memory} devuelve
\cod{process\_physical\_memory\_low} = \cod{false} y
\cod{process\_virtual\_memory\_low} = \cod{false}. No hay presión de memoria
externa sobre el proceso de SQL Server. Visto desde dentro de la máquina
virtual, subir el tope es seguro. Y el nivel de hipervisor lo confirma
—\,sin \emph{ballooning} ni intercambio, con los 128\,GB concedidos\,—, como
se detalla más adelante.

\subsection{El disco queda descartado}

La latencia real por fichero (\cod{sys.dm\_io\_virtual\_file\_stats},
acumulada desde el arranque hace 27 días) sale limpia. Los ocho ficheros con
peor latencia de lectura de toda la instancia —entre los que tienen más de
10.000 lecturas— son \textbf{todos ficheros de log}, y ninguno pasa de 9\,ms:

\begin{table}[h]
\centering
\small
\begin{tabular}{llrrr}
\toprule
\textbf{Base de datos} & \textbf{Fichero} & \textbf{ms/lectura} &
\textbf{ms/escritura} & \textbf{Escrituras} \\
\midrule
ContaAlaplana & IN3ContaTpl\_Log & 9 & 1 & — \\
D365Eventos   & D365Eventos\_log & 9 & 1 & — \\
ContaVitaCer  & log              & 8 & 1 & — \\
Informatica   & Informatica\_Log & 8 & \aviso{12} & — \\
OMS           & OMS\_log         & 8 & 1 & — \\
ContaSTN      & log              & 7 & 1 & — \\
Grupo         & In3Tplnew\_log   & 7 & 2 & 281.926.025 \\
IFR           & IFR\_log         & 7 & 2 & 363.247.671 \\
\bottomrule
\end{tabular}
\end{table}

Siete a nueve milisegundos por lectura en un fichero de log es un valor
normal, y las escrituras de log —que son las que de verdad frenan a las
transacciones— están en 1 o 2\,ms, que es bueno. El almacenamiento no está
contribuyendo al atasco.

Dos detalles menores: \cod{Informatica\_Log} escribe a 12\,ms, muy por encima
del resto, aunque con poco volumen; y los logs de \cod{Grupo} e \cod{IFR}
acumulan 9,5 y 22 millones de \emph{lecturas} respectivamente, lo que es
esperable al estar en un grupo de disponibilidad: la réplica lee el log
continuamente.

\subsection{El hipervisor queda descartado}

Tras registrar el vCenter \cod{vcsa.central.local} en DPA, la pestaña
\emph{Virtualization} correlaciona las métricas del host con las de la
instancia. Las cinco instancias monitorizadas aparecen vinculadas a su
máquina virtual; SRVPRORPDB01 corre sobre el host ESXi \cod{10.150.0.8}, en
el que no comparte sitio con ninguna otra instancia monitorizada
(PROMETEO y SRVPBI sí comparten el \cod{10.150.0.9}).

\begin{table}[h]
\centering
\small
\begin{tabular}{llrr}
\toprule
\textbf{Métrica} & \textbf{Qué mide} & \textbf{Media} & \textbf{Máximo} \\
\midrule
\cod{vmReadyTime}        & espera por CPU física   & \bien{0,09\,\%} & \bien{0,13\,\%} \\
\cod{vmCpuCoStopPercent} & desincronía entre vCPU  & \bien{0\,\%}    & \bien{0\,\%} \\
\cod{vmPercentBallooned} & memoria reclamada       & \bien{0\,\%}    & \bien{0\,\%} \\
\cod{vmBalloonSize}      & tamaño del globo        & \bien{0}        & \bien{0} \\
\cod{vmMemorySwapInRate} & intercambio del host    & \bien{0}        & \bien{0} \\
\cod{vmMemorySwapOutRate}& intercambio del host    & \bien{0}        & \bien{0} \\
\cod{vmMemoryGranted}    & memoria concedida       & 130.778\,MB     & 130.778\,MB \\
\cod{vmMemoryActive}     & conjunto activo         & 71,2\,\%        & 74,5\,\% \\
\cod{hostCpuUsagePercent}& CPU del host            & \bien{40,2\,\%} & \bien{45,5\,\%} \\
\cod{hostMemoryUsage}    & memoria del host        & \bien{39,4\,\%} & \bien{39,4\,\%} \\
\cod{vmCpuUsagePercent}  & CPU usada por la VM     & \aviso{96,4\,\%}& \aviso{114,8\,\%} \\
\bottomrule
\end{tabular}
\end{table}

\paragraph{El veredicto es limpio en todo lo que importaba.} El
\emph{CPU ready time} está en el 0,09\,\%, cuando el umbral habitual de
preocupación se sitia en torno al 5\,\% y el de problema serio por encima del
10\,\%. El \emph{co-stop} —relevante precisamente en máquinas anchas como esta,
de 16 vCPU— es exactamente cero. \textbf{La máquina virtual no está esperando
por CPU física.}

Y en memoria, las tres señales que había que descartar están a cero:
\textbf{no hay \emph{ballooning}, no hay intercambio en el host, y los 128\,GB
están concedidos en su totalidad}. El host va al 39\,\% de memoria. Esto
\textbf{confirma que los 25\,GB libres que ve el sistema operativo invitado son
reales} y respaldados por RAM física, de modo que la acción 6 —subir
\cod{max server memory}— queda desbloqueada sin reservas.

\paragraph{Un dato útil para planificar.} La VM consume prácticamente toda su
asignación de CPU (96,4\,\% de media; los picos por encima del 100\,\% que
reporta DPA se explican por el cálculo sobre frecuencia nominal con
\emph{turbo} activo). Con \emph{ready time} a cero y el host al 40\,\%, la
lectura es inequívoca: no es que a la VM le falte CPU del host, es que
\textbf{está usando los 16 vCPU que tiene}. Como el host tiene holgura, ampliar
vCPU es técnicamente viable si llegara a hacer falta —pero sigue siendo la
última acción de la lista, por el mismo motivo de antes.

\paragraph{Una falsa alarma que conviene explicar.}
\cod{hostMaxTotalDiskLatency} marca 59\,ms de media y 211\,ms de pico, cifras
que serían alarmantes. No lo son: al desglosar por dispositivo, los siete
discos de fibra (\cod{DGC Fibre Channel}) del host están todos
\textbf{por debajo de 1,5\,ms} tanto en lectura como en escritura.

\begin{table}[h]
\centering
\small
\begin{tabular}{lrr}
\toprule
\textbf{Dispositivo (host 10.150.0.8)} & \textbf{ms lectura} & \textbf{ms escritura} \\
\midrule
naa.\ldots 1367656699695a40 & \bien{0,0} & \bien{0,0} \\
naa.\ldots 5f6d65661aaa6803 & \bien{0,0} & \bien{0,0} \\
naa.\ldots 2c6b6566ef2fd7fe & \bien{0,1} & \bien{0,9} \\
naa.\ldots 8f6f6566a55ec291 & \bien{0,5} & \bien{0,0} \\
naa.\ldots 8ec6916839933029 & \bien{0,6} & \bien{0,1} \\
naa.\ldots 086f65669cc587fa & \bien{0,8} & \bien{0,7} \\
naa.\ldots a1bf9168335813e0 & \bien{0,8} & \bien{1,3} \\
\bottomrule
\end{tabular}
\end{table}

\cod{maxTotalLatency} es el máximo entre \emph{todos} los dispositivos del
host, incluidos los que no son almacén de datos. Un dispositivo local, una
unidad óptica o una ruta inactiva informando valores espurios basta para
dispararlo, y es un artefacto conocido de esa métrica en vSphere. Los
dispositivos reales están sanos, lo que además \textbf{concuerda con la
medición desde dentro del invitado} (9\,ms de peor latencia por fichero) y
confirma por dos vías independientes que el almacenamiento no es el problema.

\paragraph{Advertencia sobre la ventana.} El vCenter se ha registrado hoy, así
que solo hay \textbf{unos 70 minutos de histórico} —siete puntos con
granularidad de 10 minutos—. Los valores son inequívocos en magnitud, pero
\textbf{el pico de las 08:00 todavía no se ha observado a nivel de
hipervisor}. Conviene repetir esta comprobación dentro de unos días, con la
ventana de 7 días completa, antes de dar el punto por cerrado de forma
definitiva.

\subsection{Los registros de transacciones y el \emph{version store}}

Dos comprobaciones menores que cierran flecos, una de ellas necesaria antes
de activar RCSI.

\begin{table}[h]
\centering
\small
\begin{tabular}{lrrrr}
\toprule
\textbf{Base de datos} & \textbf{VLFs} & \textbf{Activos} &
\textbf{Tamaño medio} & \textbf{Log total} \\
\midrule
Grupo  & \aviso{1.940} & 12 & 63,9\,MB & 121,0\,GB \\
IFR    & 127 & 6 & \aviso{1.465,7\,MB} & \critico{181,8\,GB} \\
Cenusa & 97 & \aviso{82 (85\,\%)} & 110,0\,MB & 10,4\,GB \\
\bottomrule
\end{tabular}
\end{table}

Ninguno es grave, pero los tres tienen algo que señalar, y por motivos
opuestos. \cod{Grupo} tiene \textbf{1.940 VLFs}, por encima del millar que se
suele tomar como referencia: no afecta al trabajo diario, pero ralentiza el
arranque, la recuperación y las copias de log. \cod{IFR} tiene el problema
contrario —\,127 VLFs de \textbf{1,4\,GB cada uno}\,—: con fragmentos tan
grandes, un solo VLF activo inmoviliza 1,4\,GB de log. Y \cod{Cenusa} tiene
el \textbf{85\,\% de sus VLFs activos}, lo que indica que el log va apurado
frente a su tamaño.

Al margen de los VLFs, llama la atención el \textbf{tamaño absoluto de los
logs}: 181,8\,GB en \cod{IFR} y 121\,GB en \cod{Grupo}. Son coherentes con
los 363 y 281 millones de escrituras medidas antes, pero conviene verificar
que la política de copias de log es la adecuada, porque un log de ese tamaño
sugiere que en algún momento no se truncó.

\paragraph{Línea base del \emph{version store}, antes de tocar RCSI.}

\begin{table}[h]
\centering
\small
\begin{tabular}{lr}
\toprule
\textbf{TempDB, uso actual} & \textbf{Valor} \\
\midrule
\emph{Version store}   & \bien{31\,MB} \\
Objetos de usuario     & 181\,MB \\
Objetos internos       & 1.268\,MB \\
Espacio libre          & \bien{430.726\,MB (420\,GB)} \\
\bottomrule
\end{tabular}
\end{table}

El \emph{version store} ocupa hoy \textbf{31\,MB}, lo lógico con RCSI
desactivado en todas las bases. Es la cifra de partida contra la que medir
después del cambio. Y el dato relevante: hay \textbf{420\,GB libres} en
TempDB. Aunque activar RCSI en \cod{Grupo} multiplicara el \emph{version
store} por mil, seguiría sobrando sitio. \textbf{Esto reduce
considerablemente el riesgo de la acción 9}, que era la más delicada del
plan: el argumento en contra ya no puede ser el espacio de TempDB, solo la
validación funcional con la aplicación.

\subsection{Un índice que falta destaca sobre todos}

\cod{sys.dm\_db\_missing\_index\_details} —también acumulado desde el arranque
hace 27 días— ordena los candidatos por impacto estimado. El primero no está
cerca del segundo: lo supera \textbf{56 veces}.

\begin{table}[h]
\centering
\small
\begin{tabular}{llrr}
\toprule
\textbf{BD} & \textbf{Tabla} & \textbf{Impacto estimado} & \textbf{Búsquedas} \\
\midrule
Grupo & \critico{cmCargaConfirmadaLogs} & \critico{45.515.122.075} & \critico{14.889.907} \\
Grupo & entMOVIMIENTO    &   805.654.619 & 1.165.115 \\
Grupo & stnERPGDExpDoc   &   437.142.615 &   260.722 \\
Grupo & ExpExpedicion    &   422.831.787 &   141.554 \\
Grupo & entALBARAN       &   379.548.793 &     2.762 \\
Grupo & DELETED\_RECID   &   334.384.162 &    19.746 \\
Grupo & entALBARAN       &   287.714.517 &     9.724 \\
OMS   & PurchaseOrderHeader & (menor)    &       --- \\
\bottomrule
\end{tabular}
\end{table}

El candidato de cabeza es un índice sobre una sola columna,
\cod{[TipoAccion]}, en \cod{Grupo.cmCargaConfirmadaLogs}. Casi quince millones
de búsquedas en 27 días sobre una tabla sin ese índice significa recorridos
completos repetidos, y un recorrido completo es exactamente lo que dispara un
plan paralelo —el mismo paralelismo que domina las esperas de esta instancia.

Siete de los ocho candidatos están en \cod{Grupo}, que es precisamente la base
de datos donde se observaron las cadenas de bloqueo. Esto convierte la acción
4 de las propuestas, que hasta ahora era genérica, en una tarea concreta y
acotada.

\paragraph{Con una advertencia.} Las sugerencias de esta vista son
heurísticas: las produce el optimizador a partir de planes individuales, sin
tener en cuenta los índices que ya existen ni el coste de mantenerlos en
escritura. Un índice de una sola columna de baja cardinalidad como
\cod{[TipoAccion]} rara vez es el índice definitivo correcto. La cifra sirve
para saber \emph{dónde mirar}, no para aplicarla tal cual. El paso siguiente
es revisar las consultas reales que atacan esa tabla —el \emph{Tuning} de DPA
las tiene— y diseñar el índice a partir de ellas.

% =====================================================================
\newpage
\section{Consecuencias}

\paragraph{Para quien usa la aplicación.} Una sesión esperó 392 segundos
—seis minutos y medio— por un bloqueo de esquema. Esperas de 200 segundos
en clientes \cod{.Net SqlClient} desde puestos de usuario. Eso no se percibe
como lentitud, se percibe como que la aplicación se ha colgado, y el usuario
cierra y vuelve a entrar, añadiendo otra sesión a la cola.

\paragraph{Para el servidor.} 1.527 segundos de espera acumulada en una
hora. Con 202 sesiones activas, es capacidad que se paga y no se aprovecha:
los núcleos están ocupados coordinando hilos en lugar de resolviendo
consultas.

\paragraph{Para los procesos desatendidos.} Quince jobs compitiendo entre
sí. Cada uno tarda más de lo que tardaría en solitario, retiene más tiempo
sus bloqueos y empeora el cuadro justo en las horas peores.

\paragraph{El riesgo si no se toca.} Esto lleva así al menos desde agosto:
en las revisiones de entonces ya se observaron cadenas de bloqueo de más de
20 minutos y un \emph{page life expectancy} de 23 segundos en esta misma
instancia. No está degenerándose de golpe, pero tampoco se corrige solo. Un
cuadro así aguanta hasta que coincide con un pico de carga real, y entonces
el escalado de bloqueos deja de ser una molestia y pasa a ser una parada.

% =====================================================================
\section{Soluciones propuestas}

Ordenadas por relación entre esfuerzo y efecto.

\begin{longtable}{p{0.4cm}p{5.4cm}p{2.3cm}p{1.6cm}p{4.2cm}}
\toprule
\textbf{\#} & \textbf{Acción} & \textbf{Coste} & \textbf{Riesgo} &
\textbf{Qué debería pasar} \\
\midrule
\endhead
1 & \textbf{Activar \cod{AUTO\_UPDATE\_STATISTICS\_ASYNC} en \cod{Grupo}}
  & un comando, sin parada & bajo & las consultas dejan de esperar a que se
    recalculen estadísticas \\
2 & \textbf{Estadísticas con \cod{FULLSCAN} varias veces al día} sobre
    \cod{entEFECTO\_FacturasCobradas365} y
    \cod{STN\_ERP\_ANALISIS\_VENTAS\_PEDIDOS}
  & horas & bajo & estimaciones realistas $\Rightarrow$ menos planes paralelos
    innecesarios \\
3 & \textbf{Corregir los jobs saturados}: alargar el intervalo de
    \cod{D365\_SINCRONIZAR\_OMS\_DISP\_MAX} (118\,\%) y
    \cod{D365\_SYNC\_TRANSFORMAR\_DISP\_MAX} (96\,\%), u optimizarlos
  & minutos para alargar & bajo & dos procesos dejan de ocupar la máquina de
    forma continua \\
4 & \textbf{Sacar el lote nocturno de la jornada}: \cod{RAPPEL\_ZCALCULOS\-FACTURAS},
    \cod{TAREAS NOCTURNAS STN} e \cod{IndexOptimize} no deben pasar de las
    07:00
  & días, con análisis & medio & desaparece el solape con el pico de las 08:00
    y los bloqueos \cod{Sch-M} \\
5 & Escalonar los arranques coincidentes, empezando por los 8 del mismo segundo
  & horas & bajo & desaparecen las olas de paralelismo cada 10 minutos \\
6 & Subir \cod{max server memory} de 94,6 a unos 107\,GB
  & minutos, sin reinicio & bajo & sube el PLE, bajan las lecturas a disco \\
7 & Diseñar el índice de \cod{Grupo.cmCargaConfirmadaLogs} a partir de las
    consultas reales que la atacan
  & horas & bajo & 15 millones de búsquedas dejan de recorrer la tabla entera \\
8 & Comprobar en el hipervisor que esos 25\,GB están libres
  & \bien{HECHO} & --- & \bien{Confirmado}: sin \emph{ballooning} ni
    intercambio. La acción 6 es segura \\
9 & \textbf{Activar \cod{READ\_COMMITTED\_SNAPSHOT} en \cod{Grupo}}
  & ventana exclusiva + pruebas de aplicación & \aviso{medio-alto} &
    \textbf{el mayor efecto sobre el bloqueo}: los lectores dejan de bloquear
    a los escritores \\
10 & Atacar el resto de candidatos de \cod{Grupo} y las consultas del
    \emph{Tuning} de DPA
  & días & medio & menos planes paralelos y menos memoria por consulta \\
11 & Revisar transacciones abiertas en la aplicación
  & días, con desarrollo & medio & desaparecen las sesiones dormidas con bloqueos \\
12 & Valorar ampliar vCPU por encima de 16 —\,el host va al 40\,\%, hay
    holgura para hacerlo
  & compra & medio & alivia SOS\_SCHEDULER\_YIELD, pero \emph{después} de las
    anteriores \\
\bottomrule
\end{longtable}

\paragraph{Por qué este orden.} Las acciones 1 a 3 se pueden hacer esta misma
semana, no requieren parada y son reversibles. La 9 es la que más efecto
tendría sobre el síntoma que dispara la alarma —las sesiones bloqueadas— pero
es también la única que cambia la semántica de la aplicación, así que va
después de haber ganado margen con las baratas. La 12 se deja para el final a
propósito: con el 72\,\% de CPU gastándose en parte en paralelismo
innecesario, comprar núcleos antes de corregir las causas es pagar por
seguir haciendo lo mismo más rápido.

\subsection{Plan por fases}

Doce acciones en una lista no se ejecutan. Agrupadas, con una métrica de
verificación por cada una, sí.

\begin{longtable}{p{0.4cm}p{5.1cm}p{4.2cm}p{4.3cm}}
\toprule
\textbf{\#} & \textbf{Acción} & \textbf{Métrica de verificación} &
\textbf{Cómo saber que funcionó} \\
\midrule
\endhead
\multicolumn{4}{l}{\cellcolor{acento!10}\textbf{Fase 1 — esta semana.} Sin
parada, reversible, sin desarrollo.} \\
\midrule
1 & \cod{AUTO\_UPDATE\_STATISTICS\_ASYNC ON} en \cod{Grupo}
  & esperas de compilación; \cod{SQL Compilations/sec}
  & desaparecen los picos de latencia sin causa aparente \\
2 & Estadísticas \cod{FULLSCAN} 3--4 veces al día en las dos tablas de mayor
    rotación
  & \cod{modification\_counter} de esas tablas
  & se mantiene por debajo del umbral, no en 26 veces el tamaño \\
3 & Alargar el intervalo de los dos jobs saturados
  & \cod{pct\_ocupacion} de la consulta de la sección 9.2
  & por debajo del 50\,\% ambos \\
6 & \cod{max server memory} de 94,6 a 107\,GB
  & \emph{Page life expectancy}
  & de 124\,s a más de 300\,s de media \\
\midrule
\multicolumn{4}{l}{\cellcolor{acento!10}\textbf{Fase 2 — este mes.} Requiere
análisis, sin cambiar la aplicación.} \\
\midrule
4 & Sacar el lote nocturno de la jornada (RAPPEL, TAREAS NOCTURNAS,
    IndexOptimize)
  & hora de fin real de esos jobs
  & ninguno pasa de las 07:00 \\
5 & Escalonar los arranques coincidentes
  & jobs activos por segundo
  & ningún segundo con más de 2 arranques \\
7 & Índice de \cod{cmCargaConfirmadaLogs}
  & impacto estimado en \cod{dm\_db\_missing\_index\_*}
  & el candidato desaparece de la lista \\
8b & \cod{LOCK\_ESCALATION = DISABLE} en \cod{entEMPRESA},
     \cod{entORDENCARGA} y \cod{entALMACEN\_LOTE}
  & \cod{row\_lock\_wait\_in\_ms} de esas tres tablas
  & caída clara frente a las 38, 22 y 7 horas actuales \\
8c & Parametrizar \cod{INSERT StnERPArticuloLoteLinea}
  & nº de planes distintos para esa sentencia
  & de 6 a 1; baja el CPU de 31.000\,s \\
-- & Activar Query Store en \cod{Grupo}
  & ---
  & habilita demostrar la cadena causal de la fase 3 \\
\midrule
\multicolumn{4}{l}{\cellcolor{acento!10}\textbf{Fase 3 — este trimestre.}
Requiere ventana, desarrollo o compra.} \\
\midrule
9 & \cod{READ\_COMMITTED\_SNAPSHOT ON} en \cod{Grupo}
  & sesiones bloqueadas; \cod{TempDB Version Store} (hoy 31\,MB, con 420\,GB
    libres)
  & bloqueos casi a cero; el espacio de TempDB ya no es un riesgo \\
10 & Consultas del \emph{Tuning} de DPA, empezando por
     \cod{StnERPArticuloLoteLinea}
  & \% de esperas \cod{CX*}
  & por debajo del 30\,\% del total \\
10b & \textbf{Revisar \cod{ACTUALIZA\_EMPRESAS\_DIRECCIONES\ldots}} en
      \cod{Cenusa}
  & memoria concedida y lecturas lógicas por ejecución
  & bajar de 1.171\,MB y de 500.544 lecturas \\
10c & Revisar los disparadores de \cod{entMOVIMIENTO} y \cod{entALBARAN}
      en \cod{Cenusa}
  & CPU por disparo
  & reducir los 10.268\,s de \cod{tU\_entMOVIMIENTO} \\
11 & Transacciones abiertas en la aplicación
  & sesiones dormidas con transacción abierta
  & de 10 a 0 \\
12 & Ampliar vCPU
  & \cod{SOS\_SCHEDULER\_YIELD}; \emph{runnable tasks}
  & solo si tras las fases 1--2 sigue alto \\
\bottomrule
\end{longtable}

\paragraph{Una regla para las tres fases.} Un cambio a la vez, con 48 horas de
observación entre uno y otro. Si se tocan los jobs, la memoria y las
estadísticas el mismo día, el resultado será mejor o peor pero nadie sabrá por
qué, y no habrá forma de deshacer solo la parte que falló.

\subsection{Hipótesis descartadas, y con qué se descartaron}

Tan importante como lo que se recomienda es lo que \emph{no} hay que tocar.
Todo lo de esta tabla se comprobó por medición directa, no por criterio.

\begin{longtable}{p{4.3cm}p{5.2cm}p{5.2cm}}
\toprule
\textbf{Hipótesis} & \textbf{Qué se midió} & \textbf{Resultado} \\
\midrule
\endhead
Paralelismo mal configurado a nivel de instancia
  & \cod{sys.configurations}
  & MAXDOP 8 sobre 16 núcleos y \emph{cost threshold} 100, frente al 5 por
    defecto. Correcto \\
Paralelismo sobrescrito a nivel de base de datos
  & \cod{sys.database\_scoped\_configu\-rations} en \cod{Grupo}
  & MAXDOP\,=\,0 con \cod{is\_value\_default}\,=\,\cod{true}; en ese ámbito el
    0 significa «heredar la instancia». Sin sobrescritura \\
Estimador de cardinalidad antiguo
  & \cod{compatibility\_level} y \cod{LEGACY\_CARDINALITY\_ESTIMATION}
  & nivel 160 y \cod{false} en las ocho bases. Ya se usa el moderno \\
TempDB mal dimensionado
  & \cod{tempdb.sys.database\_files} y \cod{dm\_db\_file\_space\_usage}
  & 8 ficheros de datos, 99\,\% libre, log al 93\,\% libre. Correcto \\
Almacenamiento lento (invitado)
  & \cod{dm\_io\_virtual\_file\_stats} por fichero
  & peor latencia de lectura 9\,ms; escrituras de log a 1--2\,ms \\
Almacenamiento lento (host)
  & \cod{hostTotalReadLatency} y \cod{hostTotalWriteLatency} por dispositivo
  & los siete discos de fibra por debajo de 1,5\,ms \\
Proceso ajeno consumiendo CPU
  & \cod{RING\_BUFFER\_SCHEDULER\_MONITOR}
  & los procesos distintos de SQL Server, entre el 3 y el 6\,\% \\
Agotamiento del grupo de hilos
  & \cod{dm\_os\_schedulers}
  & \cod{work\_queue\_count}\,=\,0 en todos los \emph{schedulers} \\
Contención de CPU en el hipervisor
  & \cod{vmReadyTime}, \cod{vmCpuCoStopPercent}
  & 0,09\,\% y 0\,\%. La VM no espera por CPU física \\
Presión de memoria del hipervisor
  & \cod{vmPercentBallooned}, \cod{vmMemorySwap*}
  & todo a cero, 128\,GB concedidos. Los 25\,GB libres son reales \\
Confirmación síncrona del grupo de disponibilidad
  & \cod{sys.availability\_replicas}, \cod{dm\_hadr\_database\_replica\_states}
  & \cod{ASYNCHRONOUS\_COMMIT} en ambas réplicas, estado \cod{HEALTHY},
    retraso 0\,s. Las confirmaciones no esperan a la secundaria \\
\bottomrule
\end{longtable}

\paragraph{Dos correcciones sobre el propio análisis.} Conviene dejarlas
registradas, porque la hipótesis inicial de este informe era otra:

\begin{enumerate}\itemsep2pt
  \item La primera sospecha fue \textbf{paralelismo mal configurado}
        (MAXDOP 0, \emph{cost threshold} 5). Resultó falsa: ambos valores
        estaban correctos. Lo que la sustituyó —paralelismo \emph{innecesario}
        por estimaciones erróneas— explica los mismos síntomas, pero el
        remedio es completamente distinto: estadísticas, no configuración.
  \item El panel marcó el log de TempDB al \textbf{0\,\% libre}. Medido
        directamente, estaba al 93\,\%. Fue una falsa alarma de esa métrica, y
        conviene no actuar sobre ella sin confirmar.
  \item En la tabla de esperas se interpretó
        \cod{HADR\_NOTIFICATION\_DEQUEUE} (2,8\,\%) como carga del grupo de
        disponibilidad. Es una \textbf{espera inactiva} de hilos de fondo y no
        debe computarse. La sospecha derivada —que el AG estuviera en
        confirmación síncrona— se comprobó y también resultó falsa.
\end{enumerate}

\paragraph{Qué dicen estas correcciones del método.} Tres hipótesis iniciales
descartadas por medición y sustituidas por otras. Es el funcionamiento
normal de un diagnóstico, pero conviene que quede escrito: \textbf{las
recomendaciones de este informe valen lo que valen las mediciones que las
sostienen}, no la plausibilidad del razonamiento. Por eso cada acción de la
sección anterior lleva su consulta de verificación.

\subsection{Cómo medir si ha funcionado}

\begin{table}[h]
\centering
\small
\begin{tabular}{lrr}
\toprule
\textbf{Indicador} & \textbf{Hoy} & \textbf{Objetivo razonable} \\
\midrule
\% de muestras con bloqueo (7 días) & 88\,\% & por debajo del 30\,\% \\
Media de sesiones bloqueadas        & 4,9    & por debajo de 1 \\
Pico a las 8\,h                     & 9,0    & por debajo de 3 \\
Page life expectancy medio          & 124\,s & por encima de 300\,s \\
\bottomrule
\end{tabular}
\end{table}

\subsection*{Lo que no se recomienda}

\paragraph{Matar sesiones de forma rutinaria.} Resuelve el síntoma de ese
momento y oculta el problema. Solo tiene sentido como medida puntual ante
una cadena que lleve varios minutos parada.

\paragraph{Tocar todo a la vez.} Si se cambian simultáneamente los jobs, la
memoria y las consultas, no se sabrá cuál funcionó ni cuál empeoró algo.

% =====================================================================
\newpage
\section{Consultas de diagnóstico}

Todas de solo lectura. Se ejecutan sobre SRVPRORPDB01 y reproducen punto por
punto las cifras de este informe.

\subsection{Configuración de paralelismo y memoria}

\begin{lstlisting}[style=sql]
SELECT name, value_in_use, description
FROM   sys.configurations
WHERE  name IN ('max degree of parallelism',
                'cost threshold for parallelism',
                'max server memory (MB)',
                'min server memory (MB)');
\end{lstlisting}

\subsection{Núcleos y nodos NUMA}

Determina a qué valor debe acotarse MAXDOP.

\begin{lstlisting}[style=sql]
SELECT cpu_count AS nucleos_logicos,
       (SELECT COUNT(DISTINCT parent_node_id)
          FROM sys.dm_os_schedulers
         WHERE status = 'VISIBLE ONLINE') AS nodos_numa
FROM   sys.dm_os_sys_info;
\end{lstlisting}

\subsection{Esperas acumuladas}

El panel ve una muestra; esto ve el acumulado desde el último reinicio.

\begin{lstlisting}[style=sql]
SELECT TOP 15
       wait_type,
       wait_time_ms / 1000 AS espera_total_s,
       waiting_tasks_count,
       CAST(100.0 * wait_time_ms / SUM(wait_time_ms) OVER ()
            AS DECIMAL(5,2)) AS pct
FROM   sys.dm_os_wait_stats
WHERE  wait_time_ms > 0
  AND  wait_type NOT IN ('CLR_SEMAPHORE','LAZYWRITER_SLEEP','RESOURCE_QUEUE',
                         'SLEEP_TASK','SLEEP_SYSTEMTASK','WAITFOR',
                         'LOGMGR_QUEUE','CHECKPOINT_QUEUE','XE_TIMER_EVENT',
                         'REQUEST_FOR_DEADLOCK_SEARCH','BROKER_TO_FLUSH',
                         'DIRTY_PAGE_POLL','HADR_WORK_QUEUE',
                         'SP_SERVER_DIAGNOSTICS_SLEEP','QDS_ASYNC_QUEUE')
ORDER BY wait_time_ms DESC;
\end{lstlisting}

\subsection{Cadena de bloqueo en vivo}

\begin{lstlisting}[style=sql]
SELECT r.session_id,
       r.blocking_session_id,
       r.wait_type,
       r.wait_time / 1000.0 AS espera_s,
       r.command,
       DB_NAME(r.database_id) AS bd,
       s.login_name, s.host_name, s.program_name,
       t.text AS sql_actual
FROM   sys.dm_exec_requests r
JOIN   sys.dm_exec_sessions s ON s.session_id = r.session_id
OUTER APPLY sys.dm_exec_sql_text(r.sql_handle) t
WHERE  r.blocking_session_id <> 0
   OR  r.session_id IN (SELECT blocking_session_id
                          FROM sys.dm_exec_requests
                         WHERE blocking_session_id <> 0)
ORDER BY r.wait_time DESC;
\end{lstlisting}

\subsection{Sesiones dormidas reteniendo transacciones}

\begin{lstlisting}[style=sql]
SELECT s.session_id, s.login_name, s.host_name,
       s.program_name, s.status,
       DATEDIFF(SECOND, t.transaction_begin_time, SYSDATETIME()) AS abierta_s
FROM   sys.dm_tran_session_transactions st
JOIN   sys.dm_tran_active_transactions  t ON t.transaction_id = st.transaction_id
JOIN   sys.dm_exec_sessions             s ON s.session_id     = st.session_id
WHERE  s.status = 'sleeping'
ORDER BY t.transaction_begin_time;
\end{lstlisting}

\subsection{Jobs en ejecución, por nombre}

DPA muestra los jobs por su identificador hexadecimal; esto les pone nombre.

\begin{lstlisting}[style=sql]
SELECT j.name AS job, ja.start_execution_date
FROM   msdb.dbo.sysjobactivity ja
JOIN   msdb.dbo.sysjobs j ON j.job_id = ja.job_id
WHERE  ja.start_execution_date IS NOT NULL
  AND  ja.stop_execution_date  IS NULL
  AND  ja.session_id = (SELECT MAX(session_id) FROM msdb.dbo.syssessions)
ORDER BY ja.start_execution_date;
\end{lstlisting}

\subsection{Estado de TempDB}

\begin{lstlisting}[style=sql]
SELECT name, type_desc, size/128 AS tamano_mb, growth, is_percent_growth
FROM   tempdb.sys.database_files
ORDER BY type_desc, file_id;

SELECT SUM(total_page_count)*8/1024        AS tempdb_total_mb,
       SUM(unallocated_extent_page_count)*8/1024 AS tempdb_libre_mb
FROM   tempdb.sys.dm_db_file_space_usage;

SELECT total_log_size_in_bytes/1048576 AS log_total_mb,
       used_log_space_in_bytes/1048576 AS log_usado_mb,
       CAST(used_log_space_in_percent AS DECIMAL(5,2)) AS log_usado_pct
FROM   tempdb.sys.dm_db_log_space_usage;
\end{lstlisting}

\subsection{Memoria y page life expectancy}

\begin{lstlisting}[style=sql]
SELECT (SELECT cpu_count FROM sys.dm_os_sys_info) AS nucleos,
       (SELECT MAX(cntr_value)
          FROM sys.dm_os_performance_counters
         WHERE counter_name LIKE 'Page life expectancy%'
           AND object_name  LIKE '%Buffer Manager%')      AS ple_s,
       (SELECT MAX(cntr_value)/1024
          FROM sys.dm_os_performance_counters
         WHERE counter_name LIKE 'Total Server Memory (KB)%') AS sql_usa_mb,
       total_physical_memory_kb/1024     AS ram_fisica_mb,
       available_physical_memory_kb/1024 AS ram_libre_mb
FROM   sys.dm_os_sys_memory;
\end{lstlisting}

\subsection{Presión de memoria vista por el sistema operativo}

\begin{lstlisting}[style=sql]
SELECT total_physical_memory_kb/1024     AS ram_mb,
       available_physical_memory_kb/1024 AS libre_mb,
       system_memory_state_desc
FROM   sys.dm_os_sys_memory;

SELECT physical_memory_in_use_kb/1024 AS en_uso_mb,
       process_physical_memory_low,
       process_virtual_memory_low
FROM   sys.dm_os_process_memory;
\end{lstlisting}

\subsection{Latencia real por fichero}

\begin{lstlisting}[style=sql]
SELECT TOP 8
       DB_NAME(vfs.database_id)                      AS base_datos,
       mf.name                                       AS fichero,
       mf.type_desc,
       vfs.num_of_reads,
       vfs.io_stall_read_ms  / NULLIF(vfs.num_of_reads , 0) AS ms_por_lectura,
       vfs.num_of_writes,
       vfs.io_stall_write_ms / NULLIF(vfs.num_of_writes, 0) AS ms_por_escritura
FROM   sys.dm_io_virtual_file_stats(NULL, NULL) AS vfs
JOIN   sys.master_files AS mf
       ON  mf.database_id = vfs.database_id
       AND mf.file_id     = vfs.file_id
WHERE  vfs.num_of_reads > 10000
ORDER BY ms_por_lectura DESC;
\end{lstlisting}

\subsection{Índices que faltan, por impacto estimado}

\begin{lstlisting}[style=sql]
SELECT TOP 8
       DB_NAME(id.database_id) AS base_datos,
       OBJECT_NAME(id.object_id, id.database_id) AS tabla,
       CAST(igs.avg_total_user_cost
            * igs.avg_user_impact
            * (igs.user_seeks + igs.user_scans) AS BIGINT) AS impacto,
       igs.user_seeks,
       id.equality_columns,
       id.inequality_columns,
       id.included_columns
FROM   sys.dm_db_missing_index_group_stats AS igs
JOIN   sys.dm_db_missing_index_groups      AS ig
       ON ig.index_group_handle = igs.group_handle
JOIN   sys.dm_db_missing_index_details     AS id
       ON id.index_handle = ig.index_handle
ORDER BY impacto DESC;
\end{lstlisting}

\paragraph{Nota sobre la ventana de estas dos vistas.} Ambas se vacían al
reiniciar el servicio. SQL Server arrancó el 2 de octubre de 2026, así que las
cifras cubren 27 días de actividad real: una muestra amplia, no una foto
instantánea.

\subsection{Configuración por base de datos}

\begin{lstlisting}[style=sql]
SELECT name AS bd, compatibility_level AS compat,
       is_read_committed_snapshot_on   AS rcsi,
       snapshot_isolation_state_desc   AS snapshot_iso,
       is_auto_update_stats_on         AS auto_upd_stats,
       is_auto_update_stats_async_on   AS auto_upd_async,
       is_query_store_on               AS query_store
FROM   sys.databases
WHERE  database_id > 4
ORDER BY name;

USE Grupo;

SELECT name AS opcion, value, value_for_secondary, is_value_default
FROM   sys.database_scoped_configurations
ORDER BY name;
\end{lstlisting}

\subsection{Carga real del SQL Agent}

\begin{lstlisting}[style=sql]
-- Cuantas planificaciones sub-diarias hay, y a que ritmo
SELECT CASE s.freq_subday_type WHEN 2 THEN 'segundos' WHEN 4 THEN 'minutos'
            WHEN 8 THEN 'horas' ELSE 'otro' END AS unidad,
       s.freq_subday_interval AS cada_n, COUNT(*) AS num_jobs
FROM   msdb.dbo.sysjobs j
JOIN   msdb.dbo.sysjobschedules js ON js.job_id = j.job_id
JOIN   msdb.dbo.sysschedules    s  ON s.schedule_id = js.schedule_id
WHERE  j.enabled = 1 AND s.enabled = 1 AND s.freq_subday_interval > 0
GROUP BY s.freq_subday_type, s.freq_subday_interval
ORDER BY unidad, cada_n;
\end{lstlisting}

\begin{lstlisting}[style=sql]
-- Jobs que tardan mas que su propio intervalo de planificacion
WITH prog AS (
  SELECT j.job_id, j.name,
         s.freq_subday_interval * CASE s.freq_subday_type
             WHEN 2 THEN 1 WHEN 4 THEN 60 WHEN 8 THEN 3600 ELSE 60 END AS intervalo_s
  FROM   msdb.dbo.sysjobs j
  JOIN   msdb.dbo.sysjobschedules js ON js.job_id = j.job_id
  JOIN   msdb.dbo.sysschedules    s  ON s.schedule_id = js.schedule_id
  WHERE  j.enabled = 1 AND s.enabled = 1 AND s.freq_subday_interval > 0
),
dur AS (
  SELECT h.job_id,
         AVG(h.run_duration/10000*3600 + (h.run_duration/100%100)*60
             + h.run_duration%100) AS media_s,
         MAX(h.run_duration/10000*3600 + (h.run_duration/100%100)*60
             + h.run_duration%100) AS max_s,
         COUNT(*) AS ejec
  FROM   msdb.dbo.sysjobhistory h
  WHERE  h.step_id = 0
    AND  h.run_date >= CONVERT(INT, CONVERT(VARCHAR(8), DATEADD(DAY,-7,GETDATE()), 112))
  GROUP BY h.job_id
)
SELECT TOP 20 p.name AS job, p.intervalo_s, d.media_s, d.max_s, d.ejec,
       CAST(100.0*d.media_s/NULLIF(p.intervalo_s,0) AS DECIMAL(8,1)) AS pct_ocupacion
FROM   prog p JOIN dur d ON d.job_id = p.job_id
WHERE  d.max_s > p.intervalo_s
ORDER BY pct_ocupacion DESC;
\end{lstlisting}

\begin{lstlisting}[style=sql]
-- Ejecuciones largas: cuando empezaron y cuando acabaron de verdad
SELECT TOP 25 j.name AS job,
       msdb.dbo.agent_datetime(h.run_date, h.run_time) AS empezo,
       h.run_duration/10000*3600 + (h.run_duration/100%100)*60
           + h.run_duration%100 AS dur_s,
       DATEADD(SECOND, h.run_duration/10000*3600 + (h.run_duration/100%100)*60
           + h.run_duration%100,
           msdb.dbo.agent_datetime(h.run_date, h.run_time)) AS acabo
FROM   msdb.dbo.sysjobhistory h
JOIN   msdb.dbo.sysjobs j ON j.job_id = h.job_id
WHERE  h.step_id = 0
  AND  h.run_date >= CONVERT(INT, CONVERT(VARCHAR(8), DATEADD(DAY,-7,GETDATE()), 112))
  AND  h.run_duration > 3000
ORDER BY dur_s DESC;
\end{lstlisting}

\subsection{Obsolescencia de las estadísticas}

\begin{lstlisting}[style=sql]
USE Grupo;

SELECT TOP 12 OBJECT_NAME(s.object_id) AS tabla, s.name AS estadistica,
       sp.last_updated, sp.rows AS filas,
       CAST(100.0*sp.rows_sampled/NULLIF(sp.rows,0) AS DECIMAL(5,1)) AS pct_muestra,
       sp.modification_counter AS modificaciones,
       DATEDIFF(DAY, sp.last_updated, GETDATE()) AS dias
FROM   sys.stats s
CROSS APPLY sys.dm_db_stats_properties(s.object_id, s.stats_id) sp
WHERE  sp.rows > 1000000
  AND  OBJECTPROPERTY(s.object_id, 'IsUserTable') = 1
ORDER BY sp.modification_counter DESC;
\end{lstlisting}

\subsection{Métricas de hipervisor (API interna de DPA)}

Estas no son consultas SQL: DPA las obtiene del vCenter y las expone en su API
interna, la misma que consume la pestaña \emph{Virtualization}. Requieren
cookie de sesión web, no el token de la API pública.

\begin{lstlisting}[style=sql]
# Instancias y su vinculacion con la maquina virtual
GET /iwc/rest/databases
    -> id, name, linkedToVm, virtualMachineName,
       virtualMachineCurrentHostName, vsphereServerName

# Catalogo de metricas en modo Virtualization (pm=V)
GET /iwc/rest/databases/1/metrics/categories
      ?startTime=...&endTime=...&pm=V&require=metrics,status

# Serie temporal de una metrica concreta
GET /iwc/rest/databases/1/metrics/data
      ?metrics=vmReadyTime&startTime=...&endTime=...
      &useResourceGranularities=true&baseline=false
    -> data[0].series[n].dataPoints[].{time,value}

# Metricas relevantes para descartar contencion del host:
#   CPU     : vmReadyTime, vmCpuCoStopPercent, hostCpuUsagePercent,
#             vmCpuUsagePercent, vmCpuUsed, vmCpuUsageMhz
#   Memoria : vmPercentBallooned, vmBalloonSize, vmMemorySwapInRate,
#             vmMemorySwapOutRate, vmMemoryGranted, vmMemoryActive,
#             hostMemoryUsage
#   Disco   : hostMaxTotalDiskLatency, hostTotalReadLatency,
#             hostTotalWriteLatency (estas dos, desglosadas por dispositivo)
\end{lstlisting}

\subsection{Presión de CPU y de hilos}

\begin{lstlisting}[style=sql]
SELECT SUM(runnable_tasks_count) AS tareas_en_cola,
       SUM(work_queue_count)     AS sin_hilo,
       SUM(current_tasks_count)  AS tareas_totales,
       SUM(active_workers_count) AS workers,
       MAX(load_factor) AS load_factor_max,
       AVG(load_factor) AS load_factor_medio
FROM   sys.dm_os_schedulers
WHERE  status = 'VISIBLE ONLINE';
\end{lstlisting}

\begin{lstlisting}[style=sql]
-- CPU de SQL Server frente a la del resto del sistema (ultimos ~256 min)
SELECT TOP 10
       DATEADD(ms, -1*((SELECT ms_ticks FROM sys.dm_os_sys_info) - y.[timestamp]),
               GETDATE()) AS hora,
       y.cpu_sql,
       100 - y.cpu_libre - y.cpu_sql AS cpu_otros,
       y.cpu_libre
FROM (
  SELECT x.[timestamp],
         x.record.value('(./Record/SchedulerMonitorEvent/SystemHealth/SystemIdle)[1]',
                        'int') AS cpu_libre,
         x.record.value('(./Record/SchedulerMonitorEvent/SystemHealth/ProcessUtilization)[1]',
                        'int') AS cpu_sql
  FROM (
    SELECT [timestamp], CONVERT(XML, record) AS record
    FROM   sys.dm_os_ring_buffers
    WHERE  ring_buffer_type = 'RING_BUFFER_SCHEDULER_MONITOR'
      AND  record LIKE '%<SystemHealth>%'
  ) AS x
) AS y
ORDER BY y.[timestamp] DESC;
\end{lstlisting}

% =====================================================================
\newpage
\section{Evidencia}

Salida literal de la herramienta de consultas de DPA contra SRVPRORPDB01, el
29 de octubre de 2026, conectada como \cod{dpaMONITOR}.

\begin{lstlisting}[style=salida]
SELECT name, value_in_use FROM sys.configurations WHERE name IN (...)
Time: 94 ms

name                             value_in_use
cost threshold for parallelism   100
max degree of parallelism        8
max server memory (MB)           96896
\end{lstlisting}

\begin{lstlisting}[style=salida]
SELECT cpu_count, numa, page life expectancy, memoria ...
Time: 47 ms

nucleos   numa   ple_s   sql_usa_mb   ram_fisica_mb   ram_libre_mb
16        2      28      96896        131071          25657
\end{lstlisting}

\begin{lstlisting}[style=salida]
SELECT TOP 6 wait_type, wait_time_ms/1000, pct FROM sys.dm_os_wait_stats ...
Time: 15 ms

wait_type                   espera_s     pct
CXCONSUMER                  45504691     27,55
CXSYNC_PORT                 31022696     18,79
CXPACKET                    20425851     12,37
SOS_SCHEDULER_YIELD         12906011      7,82
LATCH_EX                     9896832      5,99
HADR_NOTIFICATION_DEQUEUE    4645761      2,81
\end{lstlisting}

\begin{lstlisting}[style=salida]
SELECT sqlserver_start_time, dias_encendido, virtual_machine_type_desc ...
Time: 16 ms

sqlserver_start_time      dias_encendido   virtual_machine_type_desc   schedulers
2026-10-02 14:25:05.69    27               HYPERVISOR                  16
\end{lstlisting}

\begin{lstlisting}[style=salida]
SELECT name, type_desc, size/128 ... FROM tempdb.sys.database_files
Time: 63 ms

name       type_desc   tamano_mb   growth   is_percent_growth
temp2      ROWS        54738       8192     false
temp3      ROWS        54762       8192     false
temp4      ROWS        52476       8192     false
temp5      ROWS        54714       8192     false
temp6      ROWS        51923       8192     false
temp7      ROWS        54500       8192     false
temp8      ROWS        54577       8192     false
tempdev    ROWS        54558       8192     false
templog    LOG          9928       8192     false

tempdb_total_mb   tempdb_libre_mb        log_total_mb  log_usado_mb  log_usado_pct
432252            427790                 9927          697           7,03
\end{lstlisting}

\begin{lstlisting}[style=salida]
SELECT j.name, ja.start_execution_date FROM msdb.dbo.sysjobactivity ...
Time: 47 ms

job                                            start_execution_date
BLOQUEAR EMPRESAS CON IMPAGOS                  2026-10-29 11:10:02.0
CARGAR DATOS PRODUCCIONES INFORME VALIDAR...   2026-10-29 11:00:01.0
D365_SINCRONIZAR_DATOS_COMPRAS                 2026-10-29 11:00:01.0
D365_SINCRONIZAR_OMS_DISP_MAX                  2026-10-29 11:10:02.0
D365_SYNC_BORRAR_OMS                           2026-10-29 11:00:01.0
D365_SYNC_MANTENER_PEDIDO_GRUPO_DESDED365      2026-10-29 11:10:02.0
D365_SYNC_MANTENER_RESERVAS_GRUPO_DESDED365    2026-10-29 11:08:06.0
D365_SYNC_TRANSFORMAR_DISP_MAX                 2026-10-29 11:09:08.0
STN_ERP_StnCRMSOLICITUD_AprobacionAuto         2026-10-29 11:10:02.0
STN_MARCAR_LINEAS_PREPARADO_TRASLADADO         2026-10-29 11:10:02.0
STN_SINCRONIZASTOCKRUMBO                       2026-10-29 11:10:02.0
STNAUTO_AUTOMATIZAR_ALBARANEADO                2026-10-29 11:02:06.0
stnSCALE_COPIA_SHIPPING_LOAD                   2026-10-29 11:10:02.0
stnSCALE_INT_ENT_expediciones                  2026-10-29 11:10:02.0
stnSCALE_INT_ENT_PRODUCCION                    2026-10-29 11:02:36.0
\end{lstlisting}

\begin{lstlisting}[style=salida]
SELECT ram_mb, libre_mb, system_memory_state_desc FROM sys.dm_os_sys_memory
Time: 31 ms

ram_mb    libre_mb   system_memory_state_desc
131071    25657      Available physical memory is high

SELECT en_uso_mb, process_physical_memory_low, ... FROM sys.dm_os_process_memory
Time: 16 ms

en_uso_mb   process_physical_memory_low   process_virtual_memory_low
97315       false                         false
\end{lstlisting}

\begin{lstlisting}[style=salida]
SELECT TOP 8 ... FROM sys.dm_io_virtual_file_stats(NULL, NULL) ...
Time: 125 ms

base_datos      fichero            tipo   lecturas     ms/lec   escrituras     ms/esc
ContaAlaplana   IN3ContaTpl_Log    LOG      ...            9      ...              1
D365Eventos     D365Eventos_log    LOG      ...            9      ...              1
ContaVitaCer    log                LOG      ...            8      ...              1
Informatica     Informatica_Log    LOG      ...            8      ...             12
OMS             OMS_log            LOG      ...            8      ...              1
ContaSTN        log                LOG      ...            7      ...              1
Grupo           In3Tplnew_log      LOG    9557386          7   281926025           2
IFR             IFR_log            LOG   22022516          7   363247671           2
\end{lstlisting}

\begin{lstlisting}[style=salida]
SELECT TOP 8 ... FROM sys.dm_db_missing_index_group_stats ...
Time: 78 ms

base_datos  tabla                   impacto          user_seeks   equality_columns
Grupo       cmCargaConfirmadaLogs   45515122075        14889907   [TipoAccion]
Grupo       entMOVIMIENTO             805654619         1165115   [atrIdTipoMovimiento],
                                                                  [atrIdAlbaran],
                                                                  [atrIdFactoria]
Grupo       stnERPGDExpDoc            437142615          260722   [IdDocumento]
Grupo       ExpExpedicion             422831787          141554
Grupo       entALBARAN                379548793            2762   [atrIdLibro]
Grupo       DELETED_RECID             334384162           19746
Grupo       entALBARAN                287714517            9724
OMS         PurchaseOrderHea...
\end{lstlisting}

\begin{lstlisting}[style=salida]
SELECT name AS bd, compatibility_level, is_read_committed_snapshot_on ... FROM sys.databases
Time: 140 ms

bd              compat  rcsi   snapshot_iso  auto_upd_stats  query_store  auto_upd_async
ContaAlaplana   160     false  OFF           true            false        false
ContaSTN        160     false  OFF           true            false        false
ContaVitaCer    160     false  OFF           true            false        false
D365Eventos     160     false  OFF           true            true         false
Grupo           160     false  OFF           true            false        false
IFR             160     false  OFF           true            false        false
Informatica     160     false  OFF           true            false        false
OMS             160     false  OFF           true            false        false

USE Grupo; SELECT name, value, is_value_default FROM sys.database_scoped_configurations
Time: 32 ms

opcion                                     value  value_for_secondary  is_value_default
ACCELERATED_PLAN_FORCING                   true   null                 true
ASYNC_STATS_UPDATE_WAIT_AT_LOW_PRIORITY    false  null                 true
LEGACY_CARDINALITY_ESTIMATION              false  null                 true
MAXDOP                                     0      null                 true
OPTIMIZE_FOR_AD_HOC_WORKLOADS              false  null                 true
PARAMETER_SNIFFING                         true   null                 true
QUERY_OPTIMIZER_HOTFIXES                   false  null                 true
\end{lstlisting}

\begin{lstlisting}[style=salida]
SELECT ... COUNT(*) FROM msdb.dbo.sysschedules WHERE freq_subday_interval > 0
Time: 141 ms

unidad      cada_n   num_jobs       unidad     cada_n   num_jobs
segundos        10          3       minutos        10          6
segundos        30          5       minutos        15          9
segundos        75          1       minutos        18          1
segundos        90          2       minutos        30         10
segundos       100          2       minutos        60          1
minutos          1          6       horas           1         24
minutos          2          7       horas           2          8
minutos          3         10       horas           3          3
minutos          4          4       horas           4          4
minutos          5         13       horas        6-12         17
                                    otro         1-10          4
                                    ----------------------------
                                    TOTAL                    140
\end{lstlisting}

\begin{lstlisting}[style=salida]
WITH prog AS (...), dur AS (...) SELECT TOP 20 ... WHERE d.max_s > p.intervalo_s
Time: 2484 ms

job                                     intervalo_s  media_s  max_s  ejec    pct
D365_SINCRONIZAR_OMS_DISP_MAX                   120      142    671   2380  118,3
D365_SYNC_TRANSFORMAR_DISP_MAX                  180      173    689   1977   96,1
D365_ENVIO_RECEPCIONES                           60       47    913   6202   78,3
stnSCALE_COPIA_SHIPPING_LOAD                     10        7    476  34670   70,0
D365_SINCRONIZAR_DATOS_COMPRAS                  300      200    707   1215   66,7
D365_SYNC_MANTENER_RESERVAS_GRUPO_DESDED365     240       98    667   2532   40,8
D365_SYNC_CREAR_PEDIDO_GRUPO_DESDED365          240       97   1737   2600   40,4
STN_MARCAR_LINEAS_PREPARADO_TRASLADADO          600      230   2336   1043   38,3
STN_SINCRONIZASTOCKRUMBO                        300      114    718   1433   38,0
stnSCALE_720_DATOS_PESADA_ACTUALIZA_MUELLE       10        3     80  29830   30,0
D365_SYNC_CREAR_OC_GRUPO_DESDE_D365             240       68    641   2676   28,3
stnSCALE_INT_ENT_PRODUCCION                      30        5   1781  19432   16,7
Interface APP WEB                               120       16   4590   5346   13,3
STNAUTO_AUTOMATIZAR_ALBARANEADO                  60        8   1161  10487   13,3
\end{lstlisting}

\begin{lstlisting}[style=salida]
SELECT TOP 20 j.name, AVG(run_duration), MAX(run_duration) ... (7 dias)
Time: 1158 ms

job                                        ejec  dur_media_s  dur_max_s  fallos
IndexOptimize - USER_DATABASES - Weekend      1        50360      50360       0
RAPPEL_ZCALCULOSFACTURAS                      7         8890      41020       0
TAREAS NOCTURNAS STN                          8         8801      29672       0
STN_MANTENER_ARTICULOS_CONSUMOS               7         8217      11323       0
RECOSTRUYE STOCK DOMINGO                      1         7079       7079       0
BAJA_RIESGOS                                  8         6766       8129       0
STN_ERP_CARGAR_ANALISIS_VENTAS_HISTORICO      8         3121       3960       0
STN_AVISO_DISPONIBLE_FASE                    32         2609       6065       0
IN3 RUMBO INVENTARIO                          7         2259       3136       0
DatabaseBackup - GRUPO_DATABASES - FULL      15         2202       2592       0
CREAR STN_QRY                                 1         1833       1833       0
DatabaseBackup -Grupo-FULL                    7         1624       1797       0
MC_Comprobar que imagenes tienen bien la ruta 8         1571       2114       0
DBA_CheckDB                                   1         1535       1535       0
Crear_STNCtl                                  6         1453       1575       0
D365_SYNC_BORRAR_OMS                         96         1133       1466       2
ACTUALIZAR FOTO CARTERA ARTICULO              6         1072       5067       0
AVISO STOCK MINIMO                            3         1063       1478       0
x_INFO VARIACION PRECIOS                      5         1119       1396       0
STN_CREAR_PERMISOS_VISUALIZACION_USUARIO      7         1028       1213       0
\end{lstlisting}

\begin{lstlisting}[style=salida]
USE Grupo; SELECT TOP 12 ... FROM sys.stats CROSS APPLY dm_db_stats_properties ...
Time: 22094 ms

tabla                              last_updated          filas    pct_mu  modificaciones
entEFECTO_FacturasCobradas365      2026-10-29 02:41:30  1184877    11,8      31476588
entEFECTO_FacturasCobradas365      2026-10-29 02:41:28  1184877    11,8      21329427
entEFECTO_FacturasCobradas365      2026-10-29 02:41:28  1184877    11,4      21329427
entEFECTO_FacturasCobradas365      2026-10-29 02:41:29  1184877    11,8      21329427
STN_ERP_ANALISIS_VENTAS_PEDIDOS    2026-10-29 02:48:51  6460673     1,4      12940889
STN_ERP_ANALISIS_VENTAS_PEDIDOS    2026-10-29 02:48:58  6460673     1,4      12940889
STN_ERP_ANALISIS_VENTAS_PEDIDOS    2026-10-29 02:48:59  6460673     1,4      12940889
STN_ERP_ANALISIS_VENTAS_PEDIDOS    2026-10-29 02:48:51  6460673     1,4      12940889
\end{lstlisting}

\begin{lstlisting}[style=salida]
SELECT SUM(runnable_tasks_count), SUM(work_queue_count) ... FROM sys.dm_os_schedulers
Time: 0 ms

tareas_en_cola  sin_hilo  tareas_totales  workers  load_factor_max  load_factor_medio
5               0         151             146      19               13

SELECT TOP 10 ... FROM sys.dm_os_ring_buffers (RING_BUFFER_SCHEDULER_MONITOR)
Time: 62 ms

hora                      cpu_sql   cpu_otros   cpu_libre
2026-10-29 11:35:00.503        65           4          31
2026-10-29 11:36:00.697        75           6          19
2026-10-29 11:37:00.877        83           5          12
2026-10-29 11:38:01.063        77           4          19
2026-10-29 11:39:01.253        67           3          30
2026-10-29 11:40:01.440        59           3          38
2026-10-29 11:41:01.627        80           3          17
2026-10-29 11:42:01.813        79           3          18
2026-10-29 11:43:02.003        76           3          21
2026-10-29 11:44:02.190        61           3          36
\end{lstlisting}

\begin{lstlisting}[style=salida]
GET /iwc/rest/databases   (vinculacion con vCenter)

id  name                          linkedToVm  vm             host ESXi     vcsa
3   AQUILES                       true        Aquiles        10.150.0.2    vcsa.central.local
2   PROMETEO                      true        Prometeo       10.150.0.9    vcsa.central.local
5   SRVPBI                        true        SRVPBI         10.150.0.9    vcsa.central.local
1   SRVPRORPDB01 via ATENEA2      true        SRVPRORPDB01   10.150.0.8    vcsa.central.local
4   SRVPROSCDB1 via SCALEPRODBAO  true        SRVPROSCDB1    10.150.0.27   vcsa.central.local
\end{lstlisting}

\begin{lstlisting}[style=salida]
GET /iwc/rest/databases/1/metrics/data?metrics=... (db_id=1, granularidad 600 s, 7 puntos)

metrica                   serie           media       maximo
vmReadyTime               SRVPRORPDB01      0,094       0,13
vmCpuCoStopPercent        SRVPRORPDB01      0           0
vmPercentBallooned        SRVPRORPDB01      0           0
vmBalloonSize             SRVPRORPDB01      0           0
vmMemorySwapInRate        SRVPRORPDB01      0           0
vmMemorySwapOutRate       SRVPRORPDB01      0           0
vmMemoryGranted           SRVPRORPDB01    130778      130778
vmMemoryActive            SRVPRORPDB01     71,16       74,47
vmCpuUsagePercent         SRVPRORPDB01     96,38      114,80
hostCpuUsagePercent       10.150.0.8       40,16       45,52
hostMemoryUsage           10.150.0.8       39,35       39,37
hostMaxTotalDiskLatency   10.150.0.8       59         211        <- artefacto
\end{lstlisting}

\begin{lstlisting}[style=salida]
GET .../metrics/data?metrics=hostTotalReadLatency  y  hostTotalWriteLatency
Desglose por dispositivo del host 10.150.0.8 -- descarta el artefacto anterior

dispositivo (DGC Fibre Channel Disk)        ms lectura   ms escritura
naa.6006016051bc51001367656699695a40            0,0          0,0
naa.6006016051bc51005f6d65661aaa6803            0,0          0,0
naa.6006016051bc51002c6b6566ef2fd7fe            0,1          0,9
naa.6006016051bc51008f6f6566a55ec291            0,5          0,0
naa.6006016051bc51008ec6916839933029            0,6          0,1
naa.6006016051bc5100086f65669cc587fa            0,8          0,7
naa.6006016051bc5100a1bf9168335813e0            0,8          1,3
\end{lstlisting}

% =====================================================================
\newpage
\section{Límites de este informe}

\paragraph{El histórico se diezma.} El panel limita cada serie a 720 puntos,
así que la ventana de siete días se reduce a una muestra por hora y la de 24
horas a una cada diez minutos. Los máximos citados son \textbf{cotas
inferiores}.

\paragraph{El contador de bloqueos y su lista no son simultáneos.} El
contador viene de una métrica que DPA muestrea cada minuto; la lista es una
foto en vivo tomada en otro instante. Que el contador diga 23 y la lista
muestre 3 no es un error.

\paragraph{Una métrica del panel resultó incorrecta.} El panel marcó el log
de TempDB al 0\,\% libre. Medido directamente, está al 93\,\% libre. O fue
un pico instantáneo, o un fallo de esa métrica de DPA.

\paragraph{El detalle de sesiones se refresca cada dos minutos} y se limita
a 60 filas ordenadas por tiempo de espera. Los repartos por programa, host y
base de datos salen de esa muestra, no del censo completo.

\paragraph{Las vistas acumuladas no distinguen la hora.}
\cod{dm\_os\_wait\_stats}, \cod{dm\_io\_virtual\_file\_stats} y
\cod{dm\_db\_missing\_index\_*} suman los 27 días transcurridos desde el
arranque, sin separar las horas punta de la madrugada. Sirven para
identificar qué domina, no para datar cuándo ocurre. La datación de este
informe viene del histórico del panel, no de ellas.

\paragraph{Los índices sugeridos son heurísticos.} Las cifras de impacto de
\cod{dm\_db\_missing\_index\_group\_stats} las estima el optimizador sobre
planes individuales, ignorando los índices existentes y el coste de
mantenimiento en escritura. Indican dónde mirar; no son un guion para
aplicar.

\paragraph{El nivel de hipervisor ya se ha mirado, pero con poco histórico.}
Era el único punto ciego y está cerrado en cuanto a magnitudes: \emph{ready
time} 0,09\,\%, \emph{co-stop} 0\,\%, \emph{ballooning} 0, intercambio 0. La
salvedad es la ventana: el vCenter se registró el mismo día, de modo que solo
hay unos 70 minutos de datos y \textbf{el pico de las 08:00 no se ha observado
todavía a ese nivel}. Procede repetir la comprobación con 7 días completos.

\paragraph{Sigue sin identificarse el origen del bloqueo \cod{Sch-M} de
392\,s.} Se cruzó la hora con el historial de \cod{IndexOptimize} y
\textbf{queda descartado}: ese día terminó a las 03:26:58, siete horas antes.
La hipótesis vigente —un \cod{TRUNCATE TABLE} sobre una tabla de paso— no
está comprobada, y comprobarla exige capturar el \emph{blocked process
report} o una sesión de \emph{Extended Events}.

\paragraph{\cod{Cenusa} se ha mirado por encima.} Apareció al final del
análisis y no ha recibido el mismo tratamiento que \cod{Grupo}: no se han
revisado sus estadísticas, ni sus índices ausentes, ni el contenido real de
sus 575 disparadores. Las cifras citadas son sólidas, pero el apartado es un
punto de partida, no un diagnóstico cerrado. Dado lo que ha aparecido con una
sola consulta, \textbf{merece un análisis propio}.

\paragraph{Query Store está desactivado en \cod{Grupo}} (solo está activo en
\cod{D365Eventos}). Sin él no hay histórico de planes de ejecución, que es lo
que permitiría demostrar —no inferir— que las estimaciones de cardinalidad
están provocando los planes paralelos. Activarlo sería el paso natural antes
de atacar las consultas una por una.

\subsection*{Fuentes}

Datos tomados el 29 de octubre de 2026 entre las 10:20 y las 11:25, de dos
orígenes:

\begin{enumerate}
  \item El panel de bases de datos en \cod{http://10.50.69.121:9000},
        que lee de SolarWinds DPA 2022.4. Histórico del 22 al 29 de octubre.
  \item Consultas directas a SRVPRORPDB01 a través de la herramienta de
        consultas de DPA, con la cuenta \cod{dpaMONITOR}.
  \item Métricas de hipervisor desde la pestaña \emph{Virtualization} de DPA,
        alimentada por el vCenter \cod{vcsa.central.local}, registrado el
        29 de octubre de 2026.
\end{enumerate}

% =====================================================================
\newpage
\section{Anexo: planificaciones de menos de una hora}

Hoja de trabajo para el escalonado (acción 5). Solo las planificaciones
activas con repetición en segundos o minutos; quedan fuera las horarias y las
de una vez al día.

\begin{longtable}{p{8.6cm}rp{1.3cm}p{1.6cm}p{1.6cm}}
\toprule
\textbf{Job} & \textbf{Cada} & \textbf{Unid.} & \textbf{Desde} & \textbf{Hasta} \\
\midrule
\endhead
stnSCALE\_720\_DATOS\_PESADA & 10 & seg & 00:00 & 23:59 \\
stnSCALE\_720\_DATOS\_PESADA\_ACTUALIZA\_MUELLE & 10 & seg & 05:45 & 22:00 \\
stnSCALE\_COPIA\_SHIPPING\_LOAD & 10 & seg & 00:00 & 23:59 \\
stnSCALE\_720\_DATOS\_PESADA\_ACTUALIZA\_TELEFONO & 30 & seg & 05:50 & 22:00 \\
stnSCALE\_INT\_ENT\_expediciones & 30 & seg & 05:30 & 22:30 \\
stnSCALE\_INT\_ENT\_expediciones & 30 & seg & 22:31 & 05:29 \\
stnSCALE\_INT\_ENT\_PRODUCCION & 30 & seg & 05:30 & 22:30 \\
stnSCALE\_INT\_ENT\_PRODUCCION & 30 & seg & 22:31 & 05:29 \\
stnSCALE\_INT\_ENT\_ORDENCARGA & 75 & seg & 05:30 & 22:00 \\
stnSCALE\_INT\_ENT & 90 & seg & 05:30 & 22:30 \\
stnSCALE\_INT\_ENT\_AJUSTES & 90 & seg & 05:30 & 22:30 \\
stnSCALE\_INT\_ENT\_ORDENCARGA & 100 & seg & 22:01 & 05:29 \\
stnSCALE\_INT\_ENT\_PARRILLA & 100 & seg & 00:00 & 23:59 \\
\midrule
D365\_ENVIO\_RECEPCIONES & 1 & min & 05:00 & 23:59 \\
STNAUTO\_AUTOMATIZAR\_ALBARANEADO & 1 & min & 00:00 & 23:59 \\
STNSCALE\_ERRORES\_ILSPRO & 1 & min & 00:00 & 23:59 \\
stnSCALE\_INT\_SAL & 1 & min & 05:30 & 22:30 \\
stnSCALE\_INT\_SAL\_ORDENCARGA & 1 & min & 05:30 & 22:30 \\
stnSCALE\_INT\_SAL\_PRODUCCION & 1 & min & 05:30 & 22:30 \\
D365\_SINCRONIZAR\_OMS\_DISP\_MAX & 2 & min & 05:00 & 23:59 \\
Interface APP WEB & 2 & min & 00:00 & 23:59 \\
RECALCULAR\_ALBARANES & 2 & min & 00:00 & 23:59 \\
REENVIAR ARTICULOS & 2 & min & 00:00 & 23:59 \\
REVISAR BLOQUEOS DE 7:30 a 8:30 & 2 & min & 06:00 & 09:00 \\
stn\_regenera\_op\_modificadas & 2 & min & 00:00 & 23:59 \\
stnSCALE\_EnvioTransportistas & 2 & min & 00:00 & 23:59 \\
ASPROVA\_COPIAR\_TRABAJO\_REAL & 3 & min & 00:00 & 23:59 \\
D365\_SYNC\_CREAR\_PEDIDO\_BLOQUEO\_DESDE365 & 3 & min & 00:00 & 23:59 \\
D365\_SYNC\_TRANSFORMAR\_DISP\_MAX & 3 & min & 05:00 & 23:59 \\
PREVISION\_CARGA\_CARGADATOS & 3 & min & 06:00 & 22:00 \\
STN\_OC\_ENPROCESO\_CARGA\_CARGADATOS & 3 & min & 00:00 & 23:59 \\
STNSCALE\_ERRORES & 3 & min & 00:00 & 23:59 \\
stnSCALE\_INT\_ENT & 3 & min & 22:31 & 05:29 \\
stnSCALE\_INT\_ENT\_AJUSTES & 3 & min & 22:31 & 05:29 \\
stnSCALE\_INT\_SAL & 3 & min & 22:31 & 05:29 \\
stnSCALE\_INT\_SAL\_PRODUCCION & 3 & min & 22:31 & 05:29 \\
D365\_SYNC\_CREAR\_OC\_GRUPO\_DESDE\_D365 & 4 & min & 00:00 & 23:59 \\
D365\_SYNC\_CREAR\_PEDIDO\_GRUPO\_DESDED365 & 4 & min & 00:00 & 23:59 \\
D365\_SYNC\_MANTENER\_RESERVAS\_GRUPO\_DESDED365 & 4 & min & 00:00 & 23:59 \\
stnscale\_albaranear\_mensajeria & 4 & min & 18:00 & 00:15 \\
D365\_SINCRONIZAR\_DATOS\_COMPRAS & 5 & min & 07:00 & 22:00 \\
D365\_SYNC\_MANTENER\_PEDIDO\_GRUPO\_DESDED365 & 5 & min & 00:00 & 23:59 \\
DESACCIONAR ORDENES & 5 & min & 05:00 & 23:00 \\
RECALCULAR\_PEDIDOS & 5 & min & 08:00 & 20:59 \\
REVISAR BLOQUEOS DE 7:30 a 8:30 & 5 & min & 16:55 & 17:40 \\
REVISAR BLOQUEOS DE 7:30 a 8:30 & 5 & min & 01:55 & 02:20 \\
STN\_ERP\_StnCRMSOLICITUD\_AprobacionAuto & 5 & min & 08:00 & 20:00 \\
STN\_SINCRONIZASTOCKRUMBO & 5 & min & 06:00 & 22:45 \\
stnSCALE\_720\_DATOS\_PESADA\_ACTUALIZA\_MUELLE & 5 & min & 22:00 & 05:44 \\
stnSCALE\_720\_DATOS\_PESADA\_ACTUALIZA\_TELEFONO & 5 & min & 22:00 & 05:49 \\
stnSCALE\_720\_DATOS\_PESADA\_BORRADO & 5 & min & 00:00 & 23:59 \\
stnSCALE\_EnvioMosaicoCarga & 5 & min & 00:00 & 23:59 \\
stnSCALE\_INT\_SAL\_ORDENCARGA & 5 & min & 22:31 & 05:29 \\
ACTUALIZA\_OPERACION & 10 & min & 00:00 & 23:59 \\
CENUSA - ARREGLA SERVICO OC CENUSA & 10 & min & 00:00 & 23:59 \\
Copiar SN\_PARRILLA\_LOG a Grupo & 10 & min & 00:00 & 23:59 \\
D365\_SYNC\_DATOS\_CLIENTE & 10 & min & 00:00 & 23:59 \\
STN\_MARCAR\_LINEAS\_PREPARADO\_TRASLADADO & 10 & min & 00:00 & 23:59 \\
stnscale\_albaranear\_mensajeria & 10 & min & 09:00 & 22:00 \\
ACT Estadisticas carga LEANM & 15 & min & 06:00 & 23:59 \\
ACTUALIZAR CODIGO UNICO & 15 & min & 08:00 & 16:45 \\
ASIGNACION AUTOMATICA CROSS-PARAGUAS & 15 & min & 06:00 & 21:00 \\
COMPROBAR AJUSTES STOCK & 15 & min & 00:00 & 23:59 \\
CONTROL JOBS SINFIN & 15 & min & 00:00 & 23:59 \\
D365\_SINCRONIZAR\_OMS\_DISP\_NORM & 15 & min & 07:00 & 23:59 \\
PREVISION\_CARGA\_CARGADATOS & 15 & min & 22:00 & 05:59 \\
STN\_ERP\_CARGAR\_ANALISIS\_VENTAS\_DIARIO & 15 & min & 06:00 & 23:00 \\
stnSCALE\_Revisar\_pallets\_SCALE & 15 & min & 00:00 & 23:59 \\
D365\_SYNC\_TRANSFORMAR\_DISP\_NORM & 18 & min & 05:00 & 23:59 \\
ACTUALIZAR\_FECHACOBRO\_EFECTOS\_DESDE\_CONCILIADOS & 30 & min & 08:00 & 20:00 \\
DesbloquearEmpresasSaldoNegativo & 30 & min & 08:00 & 19:00 \\
IN3 RUMBO ENVIO & 30 & min & 09:00 & 18:00 \\
stn720\_LIMPIA\_DATOS\_PESADA & 30 & min & 00:00 & 23:59 \\
stn\_crear\_stn\_bloqueo\_material\_producido & 30 & min & 00:00 & 23:59 \\
stn\_pzpedirproveedor\_CARGAR\_DATOS & 30 & min & 08:00 & 19:00 \\
STNSCALE\_ANULACIONES\_INCOMPLETAS & 30 & min & 00:00 & 23:59 \\
stnSCALE\_EnviaDatosEntregaPedido & 30 & min & 00:00 & 23:59 \\
stnSCALE\_EnviaDatosFechaEntregaPedido & 30 & min & 00:00 & 23:59 \\
stnSCALE\_EnviaDatosSuPedidoWickes & 30 & min & 00:00 & 23:59 \\
D365\_RECUPERAR\_PERDIDOS\_DATAVERSE\_CLAUDIA & 60 & min & 08:30 & 23:30 \\
\bottomrule
\end{longtable}

\noindent 80 planificaciones de menos de una hora. A estas hay que sumar 56
horarias y más de 40 diarias, hasta las 140 sub-diarias citadas en el informe.

% =====================================================================
\section{Glosario}

\begin{longtable}{p{4.4cm}p{10.4cm}}
\toprule
\textbf{Término} & \textbf{Qué significa} \\
\midrule
\endhead
CXCONSUMER, CXPACKET, EXECSYNC, CXSYNC\_PORT &
esperas de \textbf{paralelismo}: un hilo aguarda a que otros hilos de la
misma consulta terminen su parte. Normales en pequeñas cantidades;
dominantes indican que se reparte trabajo que no compensa repartir \\
\addlinespace
SOS\_SCHEDULER\_YIELD &
la sesión \textbf{espera turno de CPU}. No está bloqueada por nadie: hay más
trabajo que núcleos \\
\addlinespace
LCK\_M\_S, LCK\_M\_IS &
esperas de \textbf{bloqueo} para leer un objeto que otra sesión tiene tomado \\
\addlinespace
LCK\_M\_SCH\_S &
espera de \textbf{estabilidad de esquema}: alguien modifica la estructura de
un objeto y nadie más puede ni leerlo \\
\addlinespace
ASYNC\_NETWORK\_IO &
el servidor tiene resultados listos y \textbf{la aplicación no los lee} lo
bastante rápido \\
\addlinespace
MAXDOP &
\emph{max degree of parallelism}: cuántos núcleos puede usar una sola
consulta. A 0 significa «todos» \\
\addlinespace
Cost threshold for parallelism &
el coste estimado a partir del cual una consulta se reparte entre hilos. Por
defecto 5, un valor de 1997 \\
\addlinespace
Page life expectancy (PLE) &
segundos que una página de datos sobrevive en memoria antes de ser
desalojada. Cuanto más bajo, más vuelve el servidor al disco \\
\addlinespace
TempDB &
base de datos temporal que comparten todas las consultas del servidor para
ordenaciones, tablas temporales y versionado \\
\addlinespace
Availability Group &
conjunto de réplicas de las mismas bases. ATENEA2 es el nombre del
\emph{listener}, el punto de entrada de las aplicaciones \\
\addlinespace
SPID &
identificador de una sesión dentro del servidor \\
\addlinespace
Sesión \cod{suspended} / \cod{sleeping} &
\emph{suspended} = está esperando algo. \emph{sleeping} = no hace nada, pero
si tiene una transacción abierta sigue reteniendo sus bloqueos \\
\bottomrule
\end{longtable}

\end{document}
